\chapter{雜湊}
\label{ch:hash}

% 本章的圖共用的樣式（名稱都以 h 開頭，避免和其他章衝突）
\tikzset{
    hcell/.style={draw=black!55, minimum width=0.92cm, minimum height=0.62cm,
        inner sep=0pt, font=\ttfamily\footnotesize},
    hlnew/.style={draw=memStroke, thick, fill=memFill},
    hlprobe/.style={fill=optFill},
    hldead/.style={draw=optStroke, thick, fill=optFill},
    hbucket/.style={draw=black!55, minimum width=0.8cm, minimum height=0.6cm,
        inner sep=0pt, font=\ttfamily\footnotesize},
    hnode/.style={draw=black!55, minimum height=0.6cm, inner sep=2pt,
        font=\ttfamily\scriptsize},
    hdir/.style={draw=black!55, minimum width=0.9cm, minimum height=0.5cm,
        inner sep=0pt, font=\ttfamily\footnotesize},
    hpage/.style={draw=coreStroke, fill=coreFill, minimum width=1.9cm,
        minimum height=0.5cm, inner sep=1pt, font=\footnotesize},
    hbit/.style={draw=black!55, minimum width=0.8cm, minimum height=0.55cm,
        inner sep=0pt, font=\ttfamily\footnotesize},
    hbzero/.style={draw=optStroke, thick, fill=optStroke, text=white},
}
% 刪除標記在虛擬碼裡的寫法
\newcommand{\hashDEL}{\textsc{deleted}}

學生餐廳的刷卡機一到中午就忙不過來。 每張學生證有一個九位數的卡號，機器要在一瞬間
查出這張卡還剩多少錢。 卡號有十億種可能，開一張以卡號當索引的表太浪費；把卡號排好
再二分搜尋，或是存進第~10~章的平衡樹，每次也要走二三十步。 有沒有辦法讓每次查詢平
均只看一兩個位置，而且這個次數幾乎和全校有幾張卡無關？ 往下讀之前，先想一想。

第~1~章找重複學號時用過一個辦法：開一個以學號為索引的陣列，在上面打勾。 它快，是
因為從學號可以直接算出要看哪一格；它浪費，是因為格子數要和學號的範圍一樣多。
\textbf{雜湊}（hashing）保留「從鍵直接算出位置」的做法，但只開 $m$ 格，$m$ 和實際的
資料量差不多，再用一個\textbf{雜湊函數}（hash function）把很大的鍵空間壓到
$0, 1, \ldots, m-1$。 代價是兩個不同的鍵可能算出同一格，叫做\textbf{碰撞}
（collision）。 這一章依序講怎麼選雜湊函數、碰撞了以後怎麼存（鏈結法與開放定址）、平
均要看幾格（期望值的推導），最後講表滿了怎麼辦，以及可延伸雜湊和通用雜湊兩個觀念。 最後一節的布隆過濾器不存鍵本身，只回答
「在不在」，用一點誤判換來很大的空間節省。

雜湊表只支援插入、刪除、搜尋這三個操作，平均每次 $O(1)$。 它不維持任何順序，找最
小值、找後繼、依序列出都得把整張表掃一遍；需要這些操作時，還是要用第~9、10~章的搜
尋樹。 這一章用到《C 程式設計基礎》第~14~章的結構、第~19~章的 \texttt{malloc}、
\texttt{calloc} 和串列，以及本書第~5~章「指標的指標」的刪除寫法和第~3~章的平攤分析。

\section{直接定址表與雜湊函數}
\label{sec:hash-func}

\subsection{直接定址表}

先看鍵的範圍很小的情形。 設所有可能的鍵組成集合 $U = \{0, 1, \ldots, u - 1\}$，叫做
\textbf{鍵空間}（universe）。 \textbf{直接定址表}（direct-address table）就是一個陣
列 $T[0..u-1]$：鍵 $k$ 的資料放在 $T[k]$，沒有這個鍵時 $T[k] = \NIL$。 插入、刪除、
搜尋都只碰一格，最壞情況也是 $O(1)$。 第~1~章的 \proc{Has-Dup-Marks} 就是這種表，只
是每格只存一個「有沒有看過」。

這裡的陣列從 $0$ 開始編號，和本書虛擬碼的慣例（從 $1$ 開始）不同。 雜湊表的格子號
碼是雜湊函數算出來的餘數，從 $0$ 開始最自然，所以這一章的表一律寫成
$T[0..m-1]$，C 程式也一樣，不用換算。

直接定址表的問題是空間。 它要 $\Theta(u)$ 格，和實際存了幾個鍵無關。 九位數的卡號
$u = 10^9$，就算每格只存一個位元組也要 1 GB，而全校的卡也許只有三萬張。

\subsection{雜湊表、碰撞與載入因子}

\textbf{雜湊表}（hash table）只開 $m$ 格，$T[0..m-1]$，鍵 $k$ 放在第 $h(k)$ 格附
近。 這裡
\[
h : U \to \{0, 1, \ldots, m - 1\}
\]
是雜湊函數，$h(k)$ 叫做 $k$ 的\textbf{雜湊值}（hash value）或\textbf{起始位置}
（home slot）。 存了 $n$ 個鍵時，
\[
\alpha = \frac{n}{m}
\]
叫做\textbf{載入因子}（load factor），也就是平均每格分到幾個鍵。 之後所有的成本都用
$\alpha$ 表示。

$u > m$ 時一定有會碰撞的鍵：把整個鍵空間的 $u$ 個鍵都算一次雜湊值，由鴿籠原理，至少有
一格分到 $\lceil u/m \rceil$ 個鍵。 這還說明了一件更糟的事：不管 $h$ 怎麼選，只要 $u \ge nm$，
那一格就分到至少 $\lceil u/m \rceil \ge n$ 個鍵，也就是有 $n$ 個鍵的雜湊值全部相同。 輸入剛好是這 $n$ 個鍵時，任何雜湊表都會退化成在一
堆碰撞的鍵裡逐一比對，最壞情況是 $\Theta(n)$。 所以雜湊的分析一律談期望值，
而且要先說清楚機率從哪裡來：可能是假設鍵的雜湊值像亂數一樣均勻（第~\ref{sec:hash-chain}、
\ref{sec:hash-analysis}~節），也可能是從一族函數裡隨機挑一個（第~\ref{sec:hash-dynamic}~節
的通用雜湊）。

實用的雜湊函數要算得快，而且鍵的每一部分都會影響結果，這樣只差一兩位數的鍵（學號、
卡號常常這樣）才會被分開。 下面四種是考題最常出現的整數鍵雜湊函數。

\subsection{除法}

\textbf{除法}（division method）取餘數：
\[
h(k) = k \bmod m.
\]
算得最快，但 $m$ 不能亂選。 $m = 10^p$ 時 $h(k)$ 就是 $k$ 的最後 $p$ 位數字，前面的位
數完全沒用到；如果卡號的最後兩位是系所代碼，$m = 100$ 會讓同一系的卡全部擠在同一格。
$m = 2^p$ 時同理，只用到最低的 $p$ 個位元。 另一個陷阱是公因數。

\begin{lemma}
\label{lem:hash-gcd}
如果每個鍵都是 $d$ 的倍數，令 $g = \gcd(d, m)$，則 $k \bmod m$ 一定是 $g$ 的倍數，
所以除法最多只用到 $m/g$ 格。
\end{lemma}

\begin{proof}
寫成 $k = qm + r$，$0 \le r < m$。 $g$ 整除 $k$（因為 $g \mid d$、$d \mid k$），也整除
$m$，所以整除 $r = k - qm$。 $0$ 到 $m - 1$ 之間 $g$ 的倍數只有 $m/g$ 個。
\end{proof}

例如記憶體位址通常是 $4$ 的倍數。 鍵 $8, 20, 36, 44, 52, 64, 76, 88$ 用 $m = 12$ 時，
餘數依序是 $8, 8, 0, 8, 4, 4, 4, 4$，只用到 $0, 4, 8$ 三格；改用 $m = 13$，餘數是
$8, 7, 10, 5, 0, 12, 11, 10$，分散開了。 所以除法的 $m$ 通常選質數，而且不要
太靠近 $2$ 或 $10$ 的次方。

\subsection{乘法}

\textbf{乘法}（multiplication method）先選一個常數 $0 < A < 1$，取 $kA$ 的小數部分，
再放大到 $0$ 到 $m-1$：
\[
h(k) = \bigl\lfloor m \,( kA \bmod 1 ) \bigr\rfloor,
\]
$kA \bmod 1 = kA - \lfloor kA \rfloor$ 表示小數部分。 常用的 $A$ 是
$(\sqrt{5} - 1)/2 \approx 0.6180339887$，這個值讓連續的鍵分得很開。 例如 $m = 1000$：
\[
\begin{aligned}
5347 A &= 3304.62773\ldots, & h(5347) &= \lfloor 627.73\ldots \rfloor = 627, \\
5348 A &= 3305.24577\ldots, & h(5348) &= \lfloor 245.77\ldots \rfloor = 245.
\end{aligned}
\]
只差 $1$ 的兩個鍵落在相距很遠的兩格。 乘法對 $m$ 沒有要求，通常取 $m = 2^p$。

程式裡不用浮點數。 設一個字組有 $w$ 個位元，取整數 $s = \lfloor A \cdot 2^w \rfloor$，
$m = 2^p$。 $s \approx A \cdot 2^w$，所以 $ks \approx kA \cdot 2^w$；對 $2^w$ 取餘數，等於丟掉
$kA$ 的整數部分，剩下的低 $w$ 位元大約是 $kA$ 的小數部分乘上 $2^w$。 這個 $w$ 位元的數，
最高的 $p$ 個位元就是它除以 $2^{w-p}$ 的商，大約是小數部分乘上 $2^p = m$ 再取整數，也就是
$h(k)$：
\[
h(k) = \bigl( (k \cdot s) \bmod 2^w \bigr) \gg (w - p).
\]
在 C 裡用 \texttt{unsigned} 的乘法，溢位自動就是 $\bmod 2^w$，再右移，一共兩個運算。
用 $w = 16$、$p = 4$（$m = 16$）算 $k = 3051$：$s = \lfloor 0.6180339887 \cdot 65536
\rfloor = 40503$，$3051 \cdot 40503 = 123574653$，除以 $2^{16}$ 的餘數是
$39293 = (1001\,1001\,0111\,1101)_2$，最高 $4$ 位是 $(1001)_2 = 9$。 用實數算也一樣：
$3051A = 1885.6217\ldots$，$\lfloor 16 \times 0.6217\ldots \rfloor = 9$。

\subsection{中間平方}

\textbf{中間平方法}（mid-square method）把鍵平方，取中間的幾位。 例如四位數的鍵、
$m = 100$：$k^2$ 補成八位數，取第 $4$、$5$ 位。
\[
5347^2 = 28590409 = 285\,\underline{90}\,409, \qquad h(5347) = 90.
\]
中間的位數是由 $k$ 的每一位一起乘出來的，所以 $k$ 的每一位都會影響結果；最後幾位只
和 $k$ 的最後幾位有關，前面幾位主要由 $k$ 的前幾位決定，都不適合。 二進位的版本取
$k^2$ 中間的 $p$ 個位元，$m = 2^p$。

\subsection{摺疊}

\textbf{摺疊法}（folding）把很長的鍵切成幾段，每段的位數和 $m - 1$ 相同，再加起來，
最後只留下和 $m - 1$ 一樣多的位數（多出來的進位丟掉）。 例如九位數的卡號
$417961974$、$m = 1000$，切成 $417 \mid 961 \mid 974$：

\begin{itemize}
\item \textbf{移位摺疊}（shift folding）：直接相加，$417 + 961 + 974 = 2352$，取後三
位，$h = 352$。
\item \textbf{邊界摺疊}（boundary folding）：像摺紙一樣，每隔一段反過來，
$417 + 169 + 974 = 1560$，$h = 560$。
\end{itemize}

摺疊適合位數比 $m$ 長很多的鍵（身分證號、電話號碼）。 實務上也常先摺疊成一個整數，
再用除法或乘法。

\subsection{字串鍵}

字串 $c_0 c_1 \cdots c_{\ell-1}$ 先看成一個 $B$ 進位的數 $c_0 B^{\ell-1} + c_1
B^{\ell-2} + \cdots + c_{\ell-1}$，再對 $m$ 取餘數。 用 Horner 法（第~4~章算多項式的
方法）從左往右算，每一步都取餘數：$h \gets (h \cdot B + c_i) \bmod m$，數字就不會變大。
$B$ 選一個比字元種類多的數（例如 $131$），$m$ 選質數。 只把各字元的碼加起來不行：
\texttt{"listen"} 和 \texttt{"silent"} 用的字母一樣，和也一樣。 第~27~章的
Rabin--Karp 字串比對會用到同一個式子。

\section{鏈結法}
\label{sec:hash-chain}

處理碰撞最直接的方法是讓每一格掛一條串列：第 $j$ 格存一條串列的開頭，所有
$h(k) = j$ 的鍵都放在這條串列上。 這叫做\textbf{鏈結法}（chaining，也叫
separate chaining）。

\begin{pseudo}[label=alg:hash-chain]{鏈結法的三個操作}
\Procedure{Chain-Insert}{$T, x$}
    \State 把節點 $x$ 插在串列 $T[h(x.\mathit{key})]$ 的最前面
\EndProcedure
\Function{Chain-Search}{$T, k$}
    \State $p \gets T[h(k)]$
    \While{$p \ne \NIL$ \AND $p.\mathit{key} \ne k$}
        \State $p \gets p.\mathit{next}$
    \EndWhile
    \State \Return $p$
\EndFunction
\Procedure{Chain-Delete}{$T, x$}
    \State 把節點 $x$ 從串列 $T[h(x.\mathit{key})]$ 拿掉
\EndProcedure
\end{pseudo}

插入是 $O(1)$：插在最前面不用走到尾端。 這假設 $x$ 的鍵還不在表裡；要檢查重複，就
得先做一次搜尋。 搜尋的時間和那條串列的長度成正比。 刪除時如果用的是第~5~章的雙向
串列，給了節點 $x$ 就能 $O(1)$ 拿掉；單向串列要先找到 $x$ 的前一個節點，時間和搜尋
一樣。 下面的 C 程式用單向串列，刪除時用「指標的指標」一邊找一邊記住要改的那個
\texttt{next}。

\subsection{C 實作}
\label{sec:hash-chain-c}

章首的刷卡機用鏈結法存「卡號 $\to$ 餘額」。 為了印得下，示範只開 $m = 7$ 格，放
$8$ 張卡，$\alpha = 8/7$。 鏈結法的 $\alpha$ 可以大於 $1$，串列只是平均長一點。

\begin{cbookrecap}[第~14~章 14.3、14.5 節，第~19~章 19.3、19.9、19.10 節]
\texttt{ChainTable t = \{calloc(7, sizeof(Entry *)), 7, 0\};} 用大括號依欄位的順序
初始化結構。 \texttt{calloc} 要一塊能放 $7$ 個指標的記憶體，而且全部設成 $0$；在常
見的系統上這就是 \texttt{NULL}，所以每條串列一開始都是空的。（標準沒有保證空指標的
位元全是 $0$；要完全不靠這一點，可以改用 \texttt{malloc} 再逐格設成 \texttt{NULL}。）

\texttt{Entry **pp} 存的是「某個 \texttt{Entry *} 變數的位址」：一開始是
\texttt{\&t->bucket[b]}，往後走一步變成 \texttt{\&(*pp)->next}。 找到要刪的節點時
寫 \texttt{*pp = dead->next}，不管它是串列的第一個還是中間的，都是同一行。 這是本書
第~\ref{sec:ll-singly}~節的寫法。
\end{cbookrecap}

\begin{example}[label=ex:hash-chain]{c}{(chain\_table.c) 鏈結法的雜湊表}
#include <stdio.h>
#include <stdlib.h>

// 鏈結法的雜湊表：卡號 -> 餘額。新的鍵插在串列的最前面。
typedef struct Entry {
    long long key;
    int value;
    struct Entry *next;
} Entry;

typedef struct {
    Entry **bucket;     // bucket[0..m-1] 是 m 條串列的開頭
    int m, n;           // 格數、鍵數
} ChainTable;

// 找 key，沒找到回傳 NULL。*cmp 記錄比較了幾個鍵。
Entry *ct_find(const ChainTable *t, long long key, int *cmp)
{
    *cmp = 0;
    Entry *p = t->bucket[key % t->m];      // h(k) = k mod m
    for (; p != NULL; p = p->next) {
        (*cmp)++;
        if (p->key == key) {
            return p;
        }
    }
    return NULL;
}

// 已經有 key 就更新值，否則插在串列最前面。記憶體不夠時回傳 0。
int ct_put(ChainTable *t, long long key, int value)
{
    int cmp;
    Entry *p = ct_find(t, key, &cmp);
    if (p == NULL) {
        p = malloc(sizeof *p);
        if (p == NULL) {
            return 0;
        }
        p->key = key;
        p->next = t->bucket[key % t->m];
        t->bucket[key % t->m] = p;
        t->n++;
    }
    p->value = value;
    return 1;
}

// 刪除 key，成功回傳 1。用指標的指標，刪第一個節點不用另外處理。
int ct_delete(ChainTable *t, long long key)
{
    Entry **pp = &t->bucket[key % t->m];
    while (*pp != NULL && (*pp)->key != key) {
        pp = &(*pp)->next;
    }
    if (*pp == NULL) {
        return 0;
    }
    Entry *dead = *pp;
    *pp = dead->next;
    free(dead);
    t->n--;
    return 1;
}

void print_bucket(const ChainTable *t, int b)
{
    printf("T[%d]:", b);
    for (Entry *p = t->bucket[b]; p != NULL; p = p->next) {
        const char *arrow = (p == t->bucket[b]) ? "" : " ->";
        printf("%s %lld:%d", arrow, p->key, p->value);
    }
    printf("\n");
}

int main(void)
{
    const long long card[] = {417961974, 416402873, 417721181,
        413819238, 419364875, 410986908, 411465496, 412266090};
    const int balance[] = {350, 120, 80, 500, 45, 260, 15, 900};
    const long long ask[] = {417961974, 410986908, 414000002};
    ChainTable t = {calloc(7, sizeof(Entry *)), 7, 0};
    if (t.bucket == NULL) {
        return 1;
    }
    for (int i = 0; i < 8; i++) {
        if (!ct_put(&t, card[i], balance[i])) {
            printf("out of memory\n");
            break;
        }
    }
    for (int b = 0; b < t.m; b++) {
        print_bucket(&t, b);
    }
    for (int i = 0; i < 3; i++) {
        int cmp;
        Entry *p = ct_find(&t, ask[i], &cmp);
        int v = (p != NULL) ? p->value : -1;
        printf("get %lld: %d, cmp = %d\n", ask[i], v, cmp);
    }
    ct_put(&t, 417961974, 330);     // 已經有這張卡：只改餘額
    int first = ct_delete(&t, 417721181);
    int second = ct_delete(&t, 417721181);
    printf("delete twice: %d %d, n = %d\n", first, second, t.n);
    print_bucket(&t, 3);
    for (int b = 0; b < t.m; b++) { // 釋放：一直刪串列的第一個
        while (t.bucket[b] != NULL) {
            ct_delete(&t, t.bucket[b]->key);
        }
    }
    free(t.bucket);
    return 0;
}
\end{example}
\iosample{}{T[0]: 413819238:500 \\ T[1]: 411465496:15 \\ T[2]: \\ T[3]: 410986908:260 -> 417721181:80 -> 417961974:350 \\ T[4]: \\ T[5]: 412266090:900 -> 416402873:120 \\ T[6]: 419364875:45 \\ get 417961974: 350, cmp = 3 \\ get 410986908: 260, cmp = 1 \\ get 414000002: -1, cmp = 3 \\ delete twice: 1 0, n = 7 \\ T[3]: 410986908:260 -> 417961974:330}

前七行是放完八張卡的表，\texttt{T[3]} 那一行就是第 $3$ 格的串列。 第 $3$ 格的三張卡依放入的先後是
$417961974$、$417721181$、$410986908$，每次都插在最前面，所以順序反過來。 找
$417961974$ 要比 $3$ 次，找最後放進去的 $410986908$ 只要 $1$ 次；找不在表裡的
$414000002$ 時，它的 $h$ 也是 $3$，整條串列都要比過，比了 $3$ 次，印出 $-1$。 同
一張卡刪兩次，第二次回傳 $0$；最後一行是刪除、改餘額之後的第 $3$ 格。 程式結束前
用 \texttt{ct\_delete} 逐一刪掉每條串列的第一個節點，把記憶體還回去。 第 $2$、$4$ 格
是空的：$8$ 個鍵放進 $7$ 格，不會剛好一格一個。


\subsection{載入因子與期望長度}
\label{sec:hash-chain-analysis}

最壞的情況是 $n$ 個鍵全部落在同一格，搜尋要 $\Theta(n)$。 要說明鏈結法「平均」很
快，得先假設雜湊值怎麼分布。

\begin{definition}[簡單均勻雜湊]
\label{def:hash-suh}
每個鍵的雜湊值都等可能是 $0, 1, \ldots, m-1$ 之中的任何一個，而且不同鍵的雜湊值互
相獨立。 這個假設叫做\textbf{簡單均勻雜湊}（simple uniform hashing）。
\end{definition}

這是一個關於「雜湊函數加上輸入」的假設，不是證出來的性質。 它的意思是：雜湊函數把
實際的輸入打散得像亂數一樣。 下面的推導要用到兩個工具。 事件 $E$ 的\textbf{指示變
數}（indicator variable）$I\{E\}$ 在 $E$ 發生時是 $1$，否則是 $0$，所以
$\mathrm{E}[I\{E\}] = 1 \cdot \Pr[E] + 0 \cdot (1 - \Pr[E]) = \Pr[E]$。 期望值是線性
的：$\mathrm{E}[X + Y] = \mathrm{E}[X] + \mathrm{E}[Y]$，不需要 $X$、$Y$ 獨立（第~1~章
算插入排序的平均挪動次數用過）。 把一個要數的量寫成很多個指示變數的和，再一項一項
取期望值，是這一章最常用的手法。

\begin{theorem}
\label{thm:hash-chain-len}
在簡單均勻雜湊下，$n$ 個鍵放進 $m$ 格：
\begin{enumerate}
\item 每一條串列的期望長度都是 $\alpha = n/m$。
\item 搜尋一個不在表裡的鍵，期望比較 $\alpha$ 個鍵。
\item 搜尋一個在表裡的鍵（$n$ 個鍵被搜尋的機會一樣，新的鍵插在最前面），期望比較
\[
1 + \frac{n - 1}{2m} = 1 + \frac{\alpha}{2} - \frac{\alpha}{2n}
\]
個鍵。
\end{enumerate}
\end{theorem}

\begin{proof}
把鍵依放入的順序叫做 $x_1, x_2, \ldots, x_n$。

(1) 固定第 $j$ 格。 第 $j$ 條串列的長度是 $L_j = \sum_{i=1}^{n} I\{h(x_i) = j\}$，每
一項的期望值是 $1/m$，所以 $\mathrm{E}[L_j] = n/m$。

(2) 要找的鍵 $k$ 不在表裡，搜尋會把第 $h(k)$ 條串列整條比完。 被比到的鍵個數是
$\sum_{i=1}^{n} I\{h(x_i) = h(k)\}$。 $h(k)$ 和 $h(x_i)$ 獨立而且都均勻，所以
\[
\Pr[h(x_i) = h(k)] = \sum_{j=0}^{m-1} \Pr[h(x_i) = j] \Pr[h(k) = j]
= m \cdot \frac{1}{m^2} = \frac{1}{m},
\]
期望值是 $n/m = \alpha$。

(3) 新的鍵插在最前面，所以找 $x_i$ 時，比到的是 $x_i$ 自己，加上排在它前面、也就是
比它晚放進來而且落在同一格的鍵。 比較次數是
\[
C_i = 1 + \sum_{j = i+1}^{n} I\{h(x_j) = h(x_i)\}, \qquad
\mathrm{E}[C_i] = 1 + \frac{n - i}{m},
\]
這裡用了 (2) 的同一個計算：兩個不同的鍵落在同一格的機率是 $1/m$。 每個鍵被找的機
會是 $1/n$，所以
\[
\frac{1}{n} \sum_{i=1}^{n} \left(1 + \frac{n - i}{m}\right)
= 1 + \frac{1}{nm} \sum_{i=1}^{n} (n - i)
= 1 + \frac{1}{nm} \cdot \frac{n(n-1)}{2}
= 1 + \frac{n-1}{2m}. \qedhere
\]
\end{proof}

再加上算 $h(k)$ 和走到那一格的 $O(1)$，成功和失敗的搜尋都是期望 $\Theta(1 +
\alpha)$。 只要表的格數和鍵數成正比（$n = O(m)$，例如規定 $\alpha \le 1$），$\alpha$
就是常數，三個操作的期望時間都是 $O(1)$。 這就是章首問題的答案：查詢的次數只和
$\alpha$ 有關，和全校有幾張卡無關。

另外，(1) 說的是「隨便指定一格」的期望長度；「某個已經在表裡的鍵所在的那條串列」期望
長度是 $1 + (n-1)/m$，比 $\alpha$ 大，因為長的串列裝了比較多鍵。 期望長度是
$\alpha$ 也不代表每條都差不多長：$\alpha = 1$ 時大約有 $(1 - 1/m)^n \approx e^{-1}
\approx 37\%$ 的格子是空的（觀念題第 8 題）。

\section{開放定址}
\label{sec:hash-open}

鏈結法每個鍵要多一個指標，每次插入要 \texttt{malloc} 一個節點。 \textbf{開放定址}
（open addressing）把所有鍵直接放在表 $T[0..m-1]$ 裡，每格最多一個鍵，所以
$n \le m$、$\alpha \le 1$。 碰撞時照一個固定的順序去試別的格子，叫做\textbf{探測}
（probe）。 每個鍵 $k$ 的探測順序是
\[
h(k, 0),\ h(k, 1),\ \ldots,\ h(k, m - 1),
\]
叫做 $k$ 的\textbf{探測序列}（probe sequence），$i$ 是第幾次探測，從 $0$ 開始。
插入時沿著探測序列找第一個空格；搜尋時沿著同一個序列找，遇到 $k$ 就是找到了，遇到
空格就表示 $k$ 不在表裡：如果 $k$ 在表裡，插入它時一定會停在這個空格或更前面。 為
了讓表不滿時一定插得進去，探測序列最好是 $0, 1, \ldots, m-1$ 的一個排列，每一格都
會被試到。

\begin{pseudo}[label=alg:hash-open]{開放定址的插入與搜尋（還沒有刪除）}
\Function{Open-Insert}{$T, k$}
    \For{$i \gets 0$ \To $m - 1$}
        \State $j \gets h(k, i)$
        \If{$T[j] = \NIL$}
            \State $T[j] \gets k$
            \State \Return $j$
        \EndIf
    \EndFor
    \State \Error{表已滿}
\EndFunction
\Function{Open-Search}{$T, k$}
    \For{$i \gets 0$ \To $m - 1$}
        \State $j \gets h(k, i)$
        \If{$T[j] = k$}
            \State \Return $j$
        \ElsIf{$T[j] = \NIL$}
            \State \Return \NIL
        \EndIf
    \EndFor
    \State \Return \NIL
\EndFunction
\end{pseudo}

一次插入或搜尋的成本就是\textbf{探測次數}：看了幾格，包括讓它停下來的那一格。 下面
三種探測法的差別只在 $h(k, i)$ 怎麼算。 三種都用同一組例子：停車場發的停車證號碼，
$m = 13$，$h(k) = k \bmod 13$，依序插入
\[
5347,\ 3051,\ 6296,\ 1682,\ 2014,\ 8170,\ 2299,\ 5906,\ 7193.
\]
它們的 $h(k)$ 依序是 $4, 9, 4, 5, 12, 6, 11, 4, 4$。 放完之後
$\alpha = 9/13 \approx 0.69$。

\subsection{線性探測}

\textbf{線性探測}（linear probing）從起始位置開始一格一格往後試，到尾端繞回 $0$：
\[
h(k, i) = (h(k) + i) \bmod m.
\]
圖~\ref{fig:hash-linear} 每一列是插入一個鍵之後的表。 橘框是這次放進去的格子，粉紅底是這次探測時看過、
但已經有鍵的格子。

\begin{figure}[tbp]
\centering
\begin{tikzpicture}
\foreach \j in {0,...,12} { \node[idxlbl] at ({0.6+\j*0.92}, 0.55) {\j}; }
\node[anchor=east, font=\footnotesize] at (0.2,0.00) {放 5347，1 次};
\node[hcell] at (0.60,0.00) {};
\node[hcell] at (1.52,0.00) {};
\node[hcell] at (2.44,0.00) {};
\node[hcell] at (3.36,0.00) {};
\node[hcell, hlnew] at (4.28,0.00) {5347};
\node[hcell] at (5.20,0.00) {};
\node[hcell] at (6.12,0.00) {};
\node[hcell] at (7.04,0.00) {};
\node[hcell] at (7.96,0.00) {};
\node[hcell] at (8.88,0.00) {};
\node[hcell] at (9.80,0.00) {};
\node[hcell] at (10.72,0.00) {};
\node[hcell] at (11.64,0.00) {};
\node[anchor=east, font=\footnotesize] at (0.2,-0.78) {放 3051，1 次};
\node[hcell] at (0.60,-0.78) {};
\node[hcell] at (1.52,-0.78) {};
\node[hcell] at (2.44,-0.78) {};
\node[hcell] at (3.36,-0.78) {};
\node[hcell] at (4.28,-0.78) {5347};
\node[hcell] at (5.20,-0.78) {};
\node[hcell] at (6.12,-0.78) {};
\node[hcell] at (7.04,-0.78) {};
\node[hcell] at (7.96,-0.78) {};
\node[hcell, hlnew] at (8.88,-0.78) {3051};
\node[hcell] at (9.80,-0.78) {};
\node[hcell] at (10.72,-0.78) {};
\node[hcell] at (11.64,-0.78) {};
\node[anchor=east, font=\footnotesize] at (0.2,-1.56) {放 6296，2 次};
\node[hcell] at (0.60,-1.56) {};
\node[hcell] at (1.52,-1.56) {};
\node[hcell] at (2.44,-1.56) {};
\node[hcell] at (3.36,-1.56) {};
\node[hcell, hlprobe] at (4.28,-1.56) {5347};
\node[hcell, hlnew] at (5.20,-1.56) {6296};
\node[hcell] at (6.12,-1.56) {};
\node[hcell] at (7.04,-1.56) {};
\node[hcell] at (7.96,-1.56) {};
\node[hcell] at (8.88,-1.56) {3051};
\node[hcell] at (9.80,-1.56) {};
\node[hcell] at (10.72,-1.56) {};
\node[hcell] at (11.64,-1.56) {};
\node[anchor=east, font=\footnotesize] at (0.2,-2.34) {放 1682，2 次};
\node[hcell] at (0.60,-2.34) {};
\node[hcell] at (1.52,-2.34) {};
\node[hcell] at (2.44,-2.34) {};
\node[hcell] at (3.36,-2.34) {};
\node[hcell] at (4.28,-2.34) {5347};
\node[hcell, hlprobe] at (5.20,-2.34) {6296};
\node[hcell, hlnew] at (6.12,-2.34) {1682};
\node[hcell] at (7.04,-2.34) {};
\node[hcell] at (7.96,-2.34) {};
\node[hcell] at (8.88,-2.34) {3051};
\node[hcell] at (9.80,-2.34) {};
\node[hcell] at (10.72,-2.34) {};
\node[hcell] at (11.64,-2.34) {};
\node[anchor=east, font=\footnotesize] at (0.2,-3.12) {放 2014，1 次};
\node[hcell] at (0.60,-3.12) {};
\node[hcell] at (1.52,-3.12) {};
\node[hcell] at (2.44,-3.12) {};
\node[hcell] at (3.36,-3.12) {};
\node[hcell] at (4.28,-3.12) {5347};
\node[hcell] at (5.20,-3.12) {6296};
\node[hcell] at (6.12,-3.12) {1682};
\node[hcell] at (7.04,-3.12) {};
\node[hcell] at (7.96,-3.12) {};
\node[hcell] at (8.88,-3.12) {3051};
\node[hcell] at (9.80,-3.12) {};
\node[hcell] at (10.72,-3.12) {};
\node[hcell, hlnew] at (11.64,-3.12) {2014};
\node[anchor=east, font=\footnotesize] at (0.2,-3.90) {放 8170，2 次};
\node[hcell] at (0.60,-3.90) {};
\node[hcell] at (1.52,-3.90) {};
\node[hcell] at (2.44,-3.90) {};
\node[hcell] at (3.36,-3.90) {};
\node[hcell] at (4.28,-3.90) {5347};
\node[hcell] at (5.20,-3.90) {6296};
\node[hcell, hlprobe] at (6.12,-3.90) {1682};
\node[hcell, hlnew] at (7.04,-3.90) {8170};
\node[hcell] at (7.96,-3.90) {};
\node[hcell] at (8.88,-3.90) {3051};
\node[hcell] at (9.80,-3.90) {};
\node[hcell] at (10.72,-3.90) {};
\node[hcell] at (11.64,-3.90) {2014};
\node[anchor=east, font=\footnotesize] at (0.2,-4.68) {放 2299，1 次};
\node[hcell] at (0.60,-4.68) {};
\node[hcell] at (1.52,-4.68) {};
\node[hcell] at (2.44,-4.68) {};
\node[hcell] at (3.36,-4.68) {};
\node[hcell] at (4.28,-4.68) {5347};
\node[hcell] at (5.20,-4.68) {6296};
\node[hcell] at (6.12,-4.68) {1682};
\node[hcell] at (7.04,-4.68) {8170};
\node[hcell] at (7.96,-4.68) {};
\node[hcell] at (8.88,-4.68) {3051};
\node[hcell] at (9.80,-4.68) {};
\node[hcell, hlnew] at (10.72,-4.68) {2299};
\node[hcell] at (11.64,-4.68) {2014};
\node[anchor=east, font=\footnotesize] at (0.2,-5.46) {放 5906，5 次};
\node[hcell] at (0.60,-5.46) {};
\node[hcell] at (1.52,-5.46) {};
\node[hcell] at (2.44,-5.46) {};
\node[hcell] at (3.36,-5.46) {};
\node[hcell, hlprobe] at (4.28,-5.46) {5347};
\node[hcell, hlprobe] at (5.20,-5.46) {6296};
\node[hcell, hlprobe] at (6.12,-5.46) {1682};
\node[hcell, hlprobe] at (7.04,-5.46) {8170};
\node[hcell, hlnew] at (7.96,-5.46) {5906};
\node[hcell] at (8.88,-5.46) {3051};
\node[hcell] at (9.80,-5.46) {};
\node[hcell] at (10.72,-5.46) {2299};
\node[hcell] at (11.64,-5.46) {2014};
\node[anchor=east, font=\footnotesize] at (0.2,-6.24) {放 7193，7 次};
\node[hcell] at (0.60,-6.24) {};
\node[hcell] at (1.52,-6.24) {};
\node[hcell] at (2.44,-6.24) {};
\node[hcell] at (3.36,-6.24) {};
\node[hcell, hlprobe] at (4.28,-6.24) {5347};
\node[hcell, hlprobe] at (5.20,-6.24) {6296};
\node[hcell, hlprobe] at (6.12,-6.24) {1682};
\node[hcell, hlprobe] at (7.04,-6.24) {8170};
\node[hcell, hlprobe] at (7.96,-6.24) {5906};
\node[hcell, hlprobe] at (8.88,-6.24) {3051};
\node[hcell, hlnew] at (9.80,-6.24) {7193};
\node[hcell] at (10.72,-6.24) {2299};
\node[hcell] at (11.64,-6.24) {2014};
\end{tikzpicture}
\diagcaption{線性探測依序插入九張停車證（$m = 13$）。 左邊是插入的鍵和探測次數。}
\label{fig:hash-linear}
\end{figure}

前七個鍵都只試了一兩次。 第 $8$ 個鍵 $5906$ 的起始位置是 $4$，這時第 $4$ 到 $7$ 格已
經連成一整段，它要一路試到第 $8$ 格，共 $5$ 次。 第 $9$ 個鍵 $7193$ 也從 $4$ 開始，
這段又多了第 $8$ 格，再往後還接上本來就有鍵的第 $9$ 格，它試了 $7$ 次。 九個鍵一共
探測 $22$ 次，平均每個 $22/9 \approx 2.44$ 次。 之後搜尋這九個鍵，每個鍵的探測次數
和它當初插入時一樣，所以成功搜尋的平均也是 $2.44$ 次。

失敗的搜尋呢？ 如果要找的鍵起始位置是 $0$ 到 $3$，第一格就是空的，$1$ 次就停；起始
位置是 $4$，要把第 $4$ 到 $12$ 格看完、再看第 $0$ 格，共 $10$ 次；起始位置是 $12$，
看第 $12$、$0$ 格，共 $2$ 次。 假設起始位置是 $0$ 到 $12$ 的機會一樣，平均是
\[
\frac{4 \cdot 1 + (10 + 9 + 8 + 7 + 6 + 5 + 4 + 3 + 2)}{13} = \frac{58}{13} \approx 4.46.
\]

這種連成一段、越來越長的現象叫做\textbf{主要群聚}（primary clustering）。 原因很簡
單：一段連續的 $L$ 格後面接著一個空格，起始位置落在這 $L$ 格或那個空格本身的鍵都會
被放進那個空格，機率大約是 $(L + 1)/m$；段越長，下一個鍵越可能接在它後面，讓它更長。
兩段之間只隔一格時，那一格一填上，兩段就合成一段，上面第 $4$ 到 $12$ 格就是這樣接起
來的。

\subsection{二次探測}

\textbf{二次探測}（quadratic probing）讓跳的距離越來越大：
\[
h(k, i) = (h(k) + i^2) \bmod m,
\]
從起始位置依序跳 $0, 1, 4, 9, 16, \ldots$ 格（都取 $\bmod m$）。 起始位置不同的鍵，
就算碰上同一格，下一步也會跳到不同的地方，不會跟著同一段一路往後走，所以沒有主要群
聚。 但起始位置\textbf{相同}的鍵，探測序列整條一模一樣，例子裡 $6296$、$5906$、
$7193$ 都從 $4$ 開始，依序試 $4, 5, 8, 0$，後一個總是把前一個試過的路再走一遍。 這
叫做\textbf{次要群聚}（secondary clustering）。

二次探測有一個新問題：$i^2 \bmod m$ 不一定走得遍所有的格子。 例如 $m = 8$ 時
$i^2 \bmod 8$ 只有 $0, 1, 4$ 三種值，從起始位置出發只到得了三格，表還有空位也插不
進去（觀念題第 6 題）。 $m$ 是質數時，至少前一半是不重複的。

\begin{theorem}
\label{thm:hash-quad}
$m$ 是奇質數時，二次探測的前 $(m+1)/2$ 個位置
$h(k, 0), h(k, 1), \ldots, h(k, (m-1)/2)$ 兩兩不同。 因此只要表裡的鍵少於
$(m+1)/2$ 個（例如 $\alpha < 1/2$），插入一定成功。
\end{theorem}

\begin{proof}
假設 $0 \le i < j \le (m-1)/2$ 而 $h(k, i) = h(k, j)$。 兩邊減掉 $h(k)$，得到
$i^2 \equiv j^2 \pmod m$，也就是 $m$ 整除 $j^2 - i^2 = (j - i)(j + i)$。 $m$ 是質數，
所以 $m$ 整除 $j - i$ 或 $j + i$。 但 $1 \le j - i < m$，而且
$1 \le j + i \le m - 2 < m$，兩個都不可能被 $m$ 整除，矛盾。 前 $(m+1)/2$ 個位置兩
兩不同；表裡的鍵少於 $(m+1)/2$ 個時，這些位置裡至少有一個是空的，插入會停在那裡。
\end{proof}

這個界是緊的：$i$ 和 $m - i$ 給出同一個 $i^2 \bmod m$，所以後一半只是把前一半倒過來
再走一次，二次探測最多只看得到 $(m+1)/2$ 個不同的格子。 有兩種常見的改法讓它走遍整
張表：$m = 2^p$ 時改跳三角數 $h(k) + i(i+1)/2$（觀念題第 7 題證明它走得遍）；或是正
負交替地跳 $h(k) + 1^2, h(k) - 1^2, h(k) + 2^2, h(k) - 2^2, \ldots$，可以證明 $m$ 是
除以 $4$ 餘 $3$ 的質數時，這樣會走遍全部 $m$ 格。

\subsection{雙重雜湊}

\textbf{雙重雜湊}（double hashing）用第二個雜湊函數決定步長：
\[
h(k, i) = (h_1(k) + i \cdot h_2(k)) \bmod m.
\]
兩個鍵就算 $h_1$ 相同，$h_2$ 通常不同，探測序列在第二步就分開了，所以主要和次要群
聚都沒有。 $h_2(k)$ 不能是 $0$（會原地不動），而且要和 $m$ 互質。

\begin{theorem}
\label{thm:hash-double}
如果 $\gcd(h_2(k), m) = 1$，雙重雜湊的探測序列 $h(k, 0), \ldots, h(k, m-1)$ 是
$0, 1, \ldots, m-1$ 的一個排列。
\end{theorem}

\begin{proof}
假設 $0 \le i < j \le m - 1$ 而 $h(k, i) = h(k, j)$，則 $m$ 整除 $(j - i)\,h_2(k)$。
$m$ 和 $h_2(k)$ 互質，所以 $m$ 整除 $j - i$，但 $1 \le j - i < m$，矛盾。 $m$ 個位置
兩兩不同，就是全部的 $m$ 格。
\end{proof}

常見的兩種取法：$m$ 取質數，$h_2(k)$ 取 $1$ 到 $m - 1$ 之間的值；或是 $m = 2^p$，
$h_2(k)$ 一律取奇數。 例子裡 $m = 13$，取 $h_2(k) = 1 + (k \bmod 7)$，值在 $1$ 到
$7$ 之間，都和 $13$ 互質。

下表把三種探測法放在一起。 每一格是那個鍵插入時依序試過的格子，最後一個數就是它放
進去的位置。

\begin{center}
\begin{tabular}{rcclll}
\toprule
鍵 & $h_1$ & $h_2$ & 線性探測 & 二次探測 & 雙重雜湊 \\
\midrule
5347 & 4 & 7 & 4 & 4 & 4 \\
3051 & 9 & 7 & 9 & 9 & 9 \\
6296 & 4 & 4 & 4 5 & 4 5 & 4 8 \\
1682 & 5 & 3 & 5 6 & 5 6 & 5 \\
2014 & 12 & 6 & 12 & 12 & 12 \\
8170 & 6 & 2 & 6 7 & 6 7 & 6 \\
2299 & 11 & 4 & 11 & 11 & 11 \\
5906 & 4 & 6 & 4 5 6 7 8 & 4 5 8 & 4 10 \\
7193 & 4 & 5 & 4 5 6 7 8 9 10 & 4 5 8 0 & 4 9 1 \\
\midrule
探測總次數 & & & 22 & 17 & 13 \\
\bottomrule
\end{tabular}
\end{center}

前七個鍵三種方法幾乎一樣，差別出在最後兩個起始位置是 $4$ 的鍵。 線性探測被第 $4$
到 $10$ 格那一整段拖住；二次探測跳開了，但 $5906$、$7193$ 都重走 $4, 5, 8$；雙重雜
湊的 $7193$ 步長是 $5$，第二步就跳到第 $9$ 格，和 $6296$（步長 $4$）、$5906$（步長
$6$）走的路都不一樣。 三種方法最後的表在範例~\ref{ex:hash-open} 的輸出裡。


用「一共有幾種不同的探測序列」來比較也很清楚。 線性探測和二次探測的整條序列只由
起始位置決定，各只有 $m$ 種；雙重雜湊由 $(h_1(k), h_2(k))$ 決定，$m$ 是質數時有
$m(m-1)$ 種。 第~\ref{sec:hash-analysis}~節的理想模型允許全部 $m!$ 種排列，雙重雜湊
是三種方法裡最接近的。

\begin{examnote}
探測的追蹤題，每個鍵寫一列：$h(k)$（雙重雜湊加寫 $h_2(k)$）、依序試過的格子、最後
放的位置，最後畫出整張表。 先看題目的 $i$ 從 $0$ 還是從 $1$ 開始數，以及「探測次數」
包不包括第一次看起始位置；本章的寫法是 $i$ 從 $0$ 開始、看過幾格就算幾次。 二次探
測要寫出 $i^2 \bmod m$ 的值再加，不要心算跳過，繞回 $0$ 的地方最容易錯。
\end{examnote}

\subsection{刪除：刪除標記}
\label{sec:hash-tomb}

開放定址的刪除不能把格子改回空的。 在線性探測的表裡刪掉 $1682$（第 $6$ 格），如果
把第 $6$ 格清空，再找 $5906$：從第 $4$ 格開始，看到第 $6$ 格是空的就停下，回答「不在
表裡」。 但 $5906$ 明明在第 $8$ 格，它當初是越過第 $6$ 格才放進去的。 清空一格會把
經過它的探測序列從中間切斷。

做法是把刪掉的格子標成第三種狀態 \hashDEL，叫做\textbf{刪除標記}，也叫
墓碑（tombstone）。 搜尋時把它當成有鍵的格子，繼續往後找；插入時把它當成
空格，可以放新的鍵。 \proc{Open-Search} 一個字都不用改：\hashDEL{} 既不等於 $k$ 也不
是 \NIL，虛擬碼~\ref{alg:hash-open} 的第 11--14 行自然會跳過它。 要改的是插入和
刪除。

\begin{pseudo}[label=alg:hash-tomb]{有刪除標記的插入與刪除}
\Function{Mark-Insert}{$T, k$}
    \State $a \gets \NIL$ \Comment{第一個可以放的格子}
    \For{$i \gets 0$ \To $m - 1$}
        \State $j \gets h(k, i)$
        \If{$T[j] = k$}
            \State \Return $j$ \Comment{已經在表裡}
        \EndIf
        \If{($T[j] = \NIL$ \OR $T[j] = \hashDEL$) \AND $a = \NIL$}
            \State $a \gets j$
        \EndIf
        \If{$T[j] = \NIL$}
            \State \Break
        \EndIf
    \EndFor
    \If{$a = \NIL$}
        \State \Error{表已滿}
    \EndIf
    \State $T[a] \gets k$
    \State \Return $a$
\EndFunction
\Procedure{Mark-Delete}{$T, k$}
    \State $j \gets$ \Call{Open-Search}{$T, k$}
    \If{$j \ne \NIL$}
        \State $T[j] \gets \hashDEL$
    \EndIf
\EndProcedure
\end{pseudo}

\proc{Mark-Insert} 遇到第一個刪除標記時先記在 $a$（第 7、8 行），但不能馬上放：$k$ 可
能就在更後面，直接放會讓表裡有兩個 $k$。 所以它繼續往後找，直到遇到空格（第 9、10
行）或找到 $k$（第 5、6 行）才停，最後把 $k$ 放在第一個可以放的格子。 回到線性探測
的例子，刪除 $1682$ 之後做三件事，見圖~\ref{fig:hash-tomb}。

\begin{figure}[tbp]
\centering
\begin{tikzpicture}
\foreach \j in {0,...,12} { \node[idxlbl] at ({0.6+\j*0.92}, 0.55) {\j}; }
\node[anchor=east, font=\footnotesize] at (0.2,0.00) {刪除 1682};
\node[hcell] at (0.60,0.00) {};
\node[hcell] at (1.52,0.00) {};
\node[hcell] at (2.44,0.00) {};
\node[hcell] at (3.36,0.00) {};
\node[hcell] at (4.28,0.00) {5347};
\node[hcell, hlprobe] at (5.20,0.00) {6296};
\node[hcell, hldead] at (6.12,0.00) {$\times$};
\node[hcell] at (7.04,0.00) {8170};
\node[hcell] at (7.96,0.00) {5906};
\node[hcell] at (8.88,0.00) {3051};
\node[hcell] at (9.80,0.00) {7193};
\node[hcell] at (10.72,0.00) {2299};
\node[hcell] at (11.64,0.00) {2014};
\node[anchor=east, font=\footnotesize] at (0.2,-0.78) {找 5906，5 次};
\node[hcell] at (0.60,-0.78) {};
\node[hcell] at (1.52,-0.78) {};
\node[hcell] at (2.44,-0.78) {};
\node[hcell] at (3.36,-0.78) {};
\node[hcell, hlprobe] at (4.28,-0.78) {5347};
\node[hcell, hlprobe] at (5.20,-0.78) {6296};
\node[hcell, hlprobe] at (6.12,-0.78) {$\times$};
\node[hcell, hlprobe] at (7.04,-0.78) {8170};
\node[hcell, hlprobe] at (7.96,-0.78) {5906};
\node[hcell] at (8.88,-0.78) {3051};
\node[hcell] at (9.80,-0.78) {7193};
\node[hcell] at (10.72,-0.78) {2299};
\node[hcell] at (11.64,-0.78) {2014};
\node[anchor=east, font=\footnotesize] at (0.2,-1.56) {找 1682，9 次};
\node[hcell, hlprobe] at (0.60,-1.56) {};
\node[hcell] at (1.52,-1.56) {};
\node[hcell] at (2.44,-1.56) {};
\node[hcell] at (3.36,-1.56) {};
\node[hcell] at (4.28,-1.56) {5347};
\node[hcell, hlprobe] at (5.20,-1.56) {6296};
\node[hcell, hlprobe] at (6.12,-1.56) {$\times$};
\node[hcell, hlprobe] at (7.04,-1.56) {8170};
\node[hcell, hlprobe] at (7.96,-1.56) {5906};
\node[hcell, hlprobe] at (8.88,-1.56) {3051};
\node[hcell, hlprobe] at (9.80,-1.56) {7193};
\node[hcell, hlprobe] at (10.72,-1.56) {2299};
\node[hcell, hlprobe] at (11.64,-1.56) {2014};
\node[anchor=east, font=\footnotesize] at (0.2,-2.34) {放 4880，9 次};
\node[hcell, hlprobe] at (0.60,-2.34) {};
\node[hcell] at (1.52,-2.34) {};
\node[hcell] at (2.44,-2.34) {};
\node[hcell] at (3.36,-2.34) {};
\node[hcell] at (4.28,-2.34) {5347};
\node[hcell, hlprobe] at (5.20,-2.34) {6296};
\node[hcell, hlnew] at (6.12,-2.34) {4880};
\node[hcell, hlprobe] at (7.04,-2.34) {8170};
\node[hcell, hlprobe] at (7.96,-2.34) {5906};
\node[hcell, hlprobe] at (8.88,-2.34) {3051};
\node[hcell, hlprobe] at (9.80,-2.34) {7193};
\node[hcell, hlprobe] at (10.72,-2.34) {2299};
\node[hcell, hlprobe] at (11.64,-2.34) {2014};
\end{tikzpicture}
\diagcaption{刪除標記（$\times$）。 找 $5906$ 時越過第 $6$ 格的標記，在第 $8$ 格找到；
找已經刪掉的 $1682$ 要一路看到第 $0$ 格的空格才停；放 $4880$（$h = 5$）時記住第 $6$
格，確定後面沒有 $4880$ 之後才放回第 $6$ 格。}
\label{fig:hash-tomb}
\end{figure}

\begin{examnote}
刪除標記的題目，先看清楚題目規定插入時遇到標記怎麼做。 本章（以及多數教科書）的規
則是「繼續找到空格或找到 $k$ 為止，再放進第一個標記或空格」；有些題目為了簡化，直
接放進第一個遇到的標記。 兩種規則放的位置可能一樣，探測次數卻不同，作答時把採用的
規則寫一句在答案前面。
\end{examnote}

刪除標記讓刪除變得正確，但它不會自己消失。 插入、刪除做久了，表裡可能只剩少數幾個鍵，
卻有很多格是刪除標記；失敗的搜尋要一直走到真正的空格才停，速度由「鍵加上標記」佔的
比例決定，而不是 $\alpha$。 常見的處理是記下標記的個數 $d$，當 $(n + d)/m$ 超過某個
門檻時，把所有的鍵重新插入一張乾淨的表（第~\ref{sec:hash-dynamic}~節的再雜湊）。
刪除和放進標記格的插入都不會改變 $n + d$，只有放進真正空格的插入會讓它加一；只要
重建後的表夠空（例如 $(n + d)/m$ 不超過門檻的一半），下一次重建之前至少要再插入
$\Theta(m)$ 個鍵，重建的 $O(m)$ 平攤到這些插入上，每次 $O(1)$。 第~3~章觀念題「刪除
時先只做標記、標記多了再一起整理」用的是同一類論證。

\subsection{C 實作}
\label{sec:hash-open-c}

下面的程式把三種探測法寫在同一個函式 \texttt{probe} 裡，用 \texttt{mode} 選擇；每格
有 \texttt{EMPTY}、\texttt{FULL}、\texttt{DELETED} 三種狀態。 \texttt{search} 就
是虛擬碼的 \proc{Open-Search} 加上 \proc{Mark-Insert} 要的「第一個可以放的格子」，
\texttt{operate} 的插入、搜尋、刪除都先呼叫它，再依結果放鍵或立標記。 \texttt{main} 把這一節的九張停車證用三種方法各插一次，
線性探測再做一次刪除、搜尋和插入。 線性探測的每個操作印一行：\texttt{i}、\texttt{f}、
\texttt{d} 是插入、搜尋、刪除，接著是試過的格子、結果的格子（\texttt{-1} 表示沒找
到）和探測次數，可以和圖~\ref{fig:hash-linear}、圖~\ref{fig:hash-tomb} 逐格對照；另
外兩種方法只印最後的表，和上一小節的三列表對照。

\begin{example}[label=ex:hash-open]{c}{(open\_addr.c) 開放定址：三種探測法與刪除標記}
#include <stdio.h>
#include <stdlib.h>

// 開放定址的雜湊表（鍵是非負整數），三種探測法，刪除用標記。
enum { EMPTY, FULL, DELETED };            // 每一格的三種狀態
enum { LINEAR, QUADRATIC, DOUBLE };       // 探測法
static int trace = 1;

typedef struct {
    int *key;
    unsigned char *state;
    int m, mode;
} OpenTable;

// 第 i 次探測的格子，i 從 0 開始。h2(k) = 1 + k mod 7。
int probe(const OpenTable *t, int key, long long i)
{
    long long step = i;                           // 線性探測
    if (t->mode == QUADRATIC) {
        step = i * i;
    } else if (t->mode == DOUBLE) {
        step = i * (1 + key % 7);
    }
    return (int)((key % t->m + step) % t->m);
}

// 沿探測序列找 key，回傳它所在的格子，沒找到回傳 -1。
// *avail 設成第一個可以放 key 的格子（第一個標記，
// 或停下來的空格），沒有就是 -1。trace 是 1 時印出看過的格子。
int search(const OpenTable *t, int key, int *avail, int *probes)
{
    *avail = -1;
    for (int i = 0; i < t->m; i++) {
        int j = probe(t, key, i);
        *probes = i + 1;
        if (trace) {
            printf(" %d", j);
        }
        if (t->state[j] == FULL && t->key[j] == key) {
            return j;
        }
        if (t->state[j] != FULL && *avail < 0) {
            *avail = j;
        }
        if (t->state[j] == EMPTY) {
            return -1;
        }
    }
    return -1;
}

// op 是 'i' 插入、'f' 搜尋、'd' 刪除，每個操作都先搜尋。
void operate(OpenTable *t, char op, int key)
{
    int avail, probes;
    if (trace) {
        printf("%c %d:", op, key);
    }
    int j = search(t, key, &avail, &probes);
    if (op == 'i' && j < 0 && avail >= 0) {
        t->key[avail] = key;
        t->state[avail] = FULL;
        j = avail;
    } else if (op == 'd' && j >= 0) {
        t->state[j] = DELETED;      // 立標記，不能改回 EMPTY
    }
    if (trace) {
        printf(" -> %d, probes = %d\n", j, probes);
    }
}

void print_table(const OpenTable *t)
{
    printf("table:");
    for (int j = 0; j < t->m; j++) {
        if (t->state[j] == FULL) {
            printf(" %d", t->key[j]);
        } else {
            printf(" %s", t->state[j] == EMPTY ? "-" : "D");
        }
    }
    printf("\n");
}

int main(void)
{
    const int permit[] = {5347, 3051, 6296, 1682, 2014,
                          8170, 2299, 5906, 7193};
    const char *name[] = {"linear", "quadratic", "double"};
    for (int mode = LINEAR; mode <= DOUBLE; mode++) {
        OpenTable t = {malloc(13 * sizeof(int)), calloc(13, 1),
                       13, mode};    // calloc：全部是 EMPTY
        if (t.key != NULL && t.state != NULL) {
            trace = (mode == LINEAR);   // 只印線性探測的過程
            printf("== %s\n", name[mode]);
            for (int i = 0; i < 9; i++) {
                operate(&t, 'i', permit[i]);
            }
            print_table(&t);
            if (mode == LINEAR) {       // 刪除、搜尋、再插入
                operate(&t, 'd', 1682);
                operate(&t, 'f', 5906);
                operate(&t, 'f', 1682);
                operate(&t, 'i', 4880);
                print_table(&t);
            }
        }
        free(t.key);
        free(t.state);
    }
    return 0;
}
\end{example}
\iosample{}{== linear \\ i 5347: 4 -> 4, probes = 1 \\ i 3051: 9 -> 9, probes = 1 \\ i 6296: 4 5 -> 5, probes = 2 \\ i 1682: 5 6 -> 6, probes = 2 \\ i 2014: 12 -> 12, probes = 1 \\ i 8170: 6 7 -> 7, probes = 2 \\ i 2299: 11 -> 11, probes = 1 \\ i 5906: 4 5 6 7 8 -> 8, probes = 5 \\ i 7193: 4 5 6 7 8 9 10 -> 10, probes = 7 \\ table: - - - - 5347 6296 1682 8170 5906 3051 7193 2299 2014 \\ d 1682: 5 6 -> 6, probes = 2 \\ f 5906: 4 5 6 7 8 -> 8, probes = 5 \\ f 1682: 5 6 7 8 9 10 11 12 0 -> -1, probes = 9 \\ i 4880: 5 6 7 8 9 10 11 12 0 -> 6, probes = 9 \\ table: - - - - 5347 6296 4880 8170 5906 3051 7193 2299 2014 \\ == quadratic \\ table: 7193 - - - 5347 6296 1682 8170 5906 3051 - 2299 2014 \\ == double \\ table: - 7193 - - 5347 1682 8170 - 6296 3051 5906 2299 2014}

輸出的 \texttt{table} 一行由第 $0$ 格印到第 $12$ 格，\texttt{-} 是空格。 刪除後的
\texttt{table} 裡看不到 \texttt{D}，是因為最後插入的 $4880$ 剛好放進了第 $6$ 格的刪
除標記。 \texttt{probe} 用 \texttt{long long} 算 \texttt{i * i} 和
\texttt{i * (1 + key \% 7)}，$m$ 很大時也不會溢位。

\section{期望探測次數}
\label{sec:hash-analysis}

第~\ref{sec:hash-chain-analysis}~節算出鏈結法的期望比較次數。 這一節對開放定址做同
樣的事。 先要一個和簡單均勻雜湊對應的假設，它講的是整條探測序列。

\begin{definition}[均勻雜湊]
\label{def:hash-uh}
每個鍵的探測序列 $\langle h(k, 0), \ldots, h(k, m-1) \rangle$ 等可能是
$0, 1, \ldots, m-1$ 的 $m!$ 種排列之中的任何一種，而且不同鍵的探測序列互相獨立。
這個假設叫做\textbf{均勻雜湊}（uniform hashing）。
\end{definition}

上一節的三種方法都達不到這個假設（線性和二次探測只有 $m$ 種序列，雙重雜湊大約
$m^2$ 種）；它是一個理想模型，實際的方法離它多遠，最後用程式量。 這一節假設表裡沒有
刪除標記。

\subsection{失敗的搜尋與插入}

\begin{theorem}
\label{thm:hash-unsucc}
在均勻雜湊下，開放定址的表有 $m$ 格、$n$ 個鍵，$n < m$。 搜尋一個不在表裡的鍵，期
望探測次數是
\[
\frac{m + 1}{m - n + 1} < \frac{1}{1 - \alpha}.
\]
\end{theorem}

\begin{proof}
表裡哪 $n$ 格有鍵是由其他鍵的探測序列決定的；要找的鍵 $k$ 的探測序列和它們獨立，所
以可以先把有鍵的 $n$ 格固定下來，叫它們「滿格」，其餘 $m - n$ 格叫「空格」。 $k$ 不
在表裡，搜尋會沿著 $k$ 的探測序列一直看，直到第一個空格才停。 所以
\[
\text{探測次數} = 1 + (\text{排在第一個空格前面的滿格個數}).
\]
對每一個滿格 $s$，令 $I_s$ 是「$s$ 在 $k$ 的探測序列裡排在所有空格前面」的指示變數。
$k$ 的探測序列是均勻隨機的排列，所以 $s$ 和 $m - n$ 個空格這 $m - n + 1$ 格之間的先
後順序也是均勻的，每一格排在最前面的機率都一樣，都是 $1/(m - n + 1)$。 因此
$\mathrm{E}[I_s] = 1/(m - n + 1)$，再用期望值的線性：
\[
\mathrm{E}[\text{探測次數}] = 1 + \sum_{s} \mathrm{E}[I_s]
= 1 + \frac{n}{m - n + 1} = \frac{m + 1}{m - n + 1}.
\]
最後 $m - n + 1 > m - n$，所以這個值小於
$1 + \dfrac{n}{m - n} = \dfrac{m}{m - n} = \dfrac{1}{1 - \alpha}$。
\end{proof}

\begin{corollary}
\label{cor:hash-insert}
在均勻雜湊下，把一個新的鍵插入有 $n$ 個鍵的表，期望探測次數是
$(m + 1)/(m - n + 1)$，小於 $1/(1 - \alpha)$。
\end{corollary}

插入就是做一次失敗的搜尋，再把鍵放進停下來的那個空格，探測序列完全一樣。

$1/(1 - \alpha)$ 有一個好記的意思：每看一格，是空格的機率大約是 $1 - \alpha$，像擲一
枚正面機率 $1 - \alpha$ 的硬幣，期望要擲 $1/(1 - \alpha)$ 次才出現第一個正面。 半滿
時要 $2$ 次，$90\%$ 滿時要 $10$ 次。 表越滿，成本漲得越快。

\subsection{成功的搜尋}

\begin{theorem}
\label{thm:hash-succ}
在均勻雜湊下，表有 $m$ 格、$n$ 個鍵，$n < m$，表裡的每個鍵被搜尋的機會一樣。 成功
搜尋的期望探測次數是
\[
\frac{1}{n} \sum_{i=0}^{n-1} \frac{m+1}{m-i+1}
\;\le\; \frac{1}{\alpha} \ln \frac{1}{1 - \alpha}.
\]
\end{theorem}

\begin{proof}
把鍵依插入的順序編號。 第 $i + 1$ 個鍵插入時表裡有 $i$ 個鍵。 因為沒有刪除，格子只
會由空變滿，搜尋這個鍵時走的路和當初插入它時完全一樣：前面那些格子當時是滿的，現在
還是滿的，它停的那一格現在放的就是它。 所以搜尋它的探測次數等於插入它的探測次數，
由系理~\ref{cor:hash-insert}，期望值是 $(m + 1)/(m - i + 1)$。 每個鍵被搜尋的機率
是 $1/n$，把 $i = 0$ 到 $n - 1$ 平均就是左邊的式子。

接著估上界。 $(m + 1)(m - i) \le m(m - i + 1)$ 展開後是 $-i \le 0$，成立，所以每一項
$\frac{m+1}{m-i+1} \le \frac{m}{m-i}$。 於是
\[
\frac{1}{n} \sum_{i=0}^{n-1} \frac{m+1}{m-i+1}
\le \frac{m}{n} \sum_{i=0}^{n-1} \frac{1}{m - i}
= \frac{1}{\alpha} \sum_{j = m-n+1}^{m} \frac{1}{j}.
\]
$1/x$ 是遞減的，在區間 $[j - 1, j]$ 上 $1/x \ge 1/j$，所以
$\frac{1}{j} \le \int_{j-1}^{j} \frac{dx}{x}$。 把 $j = m - n + 1$ 到 $m$ 加起來：
\[
\sum_{j=m-n+1}^{m} \frac{1}{j} \le \int_{m-n}^{m} \frac{dx}{x}
= \ln \frac{m}{m - n} = \ln \frac{1}{1 - \alpha}. \qedhere
\]
\end{proof}

左邊的式子可以寫成 $\frac{m+1}{n}(H_{m+1} - H_{m-n+1})$，$H_t$ 是第~1~章的調和數。
用第~\ref{sec:hash-open}~節的例子驗算：$m = 13$、$n = 9$，失敗搜尋是
$14/5 = 2.8$，成功搜尋是 $\frac{14}{9}(H_{14} - H_5) \approx 1.51$。 線性探測在同一張
表上實際是 $4.46$ 和 $2.44$，雙重雜湊的成功搜尋是 $13/9 \approx 1.44$。 九個鍵太少，
這些數只能看出方向；大量的資料要用程式量。

\subsection{線性探測的公式}

線性探測有主要群聚，均勻雜湊的公式不適用。 Knuth 在 \textit{The Art of Computer
Programming} 第 3 冊 6.4 節分析了線性探測（假設起始位置 $h(k)$ 均勻而且獨立），$m$
很大、$\alpha$ 固定時，期望探測次數趨近於
\[
\text{成功：} \frac{1}{2}\left(1 + \frac{1}{1 - \alpha}\right), \qquad
\text{失敗：} \frac{1}{2}\left(1 + \frac{1}{(1 - \alpha)^2}\right).
\]
這裡只引用結果，推導要用到比本書多的組合計數，見延伸閱讀。 雙重雜湊雖然只有大約
$m^2$ 種探測序列，已經證明它在 $m$ 很大時的期望探測次數趨近於均勻雜湊的公式（Lueker
與 Molodowitch，1993）。 幾個 $\alpha$ 的公式值，下一小節的程式會和實測值一起印出來。

$\alpha = 0.5$ 時，線性探測的失敗搜尋是 $2.5$ 次，均勻雜湊是 $2$ 次，差不多；過了 $0.7$ 左右，線性探測的失敗搜尋就比均勻雜湊
貴很多，$0.9$ 時差了五倍。 所以開放定址的表通常把 $\alpha$ 控制在 $0.5$ 到 $0.75$
以下，超過就擴表。 反過來說，線性探測每次看的是相鄰的格子，在有快取的真實電腦上一
次讀進來好幾格，$\alpha$ 不高時常常比雙重雜湊快。 RAM 模型（第~1~章）數不出這種
差別。

\subsection{實測}
\label{sec:hash-measure}

下面的程式開一張 $m = 100003$ 格（質數）的表，用自己寫的線性同餘亂數產生鍵，載入因
子依序加到 $0.10$、$0.25$、$0.50$、$0.75$、$0.90$、$0.95$。 每到一個載入因子，就把
表裡每個鍵搜尋一次，算成功搜尋的平均；再搜尋 $100000$ 個一定不在表裡的鍵，算失敗搜
尋的平均。 放進表裡的鍵都是偶數、拿來做失敗搜尋的都是奇數，這樣不用另外檢查。 線性
探測和雙重雜湊（$h_2(k) = 1 + (k \bmod (m - 2))$）用同一串鍵，旁邊印出對應的公式值。

亂數用固定的種子，所以每次執行的輸出都一樣，可以逐字對照。 C 標準函式庫的
\texttt{rand} 在不同的系統上演算法不同，同一個種子會得到不同的數，這裡不用它。 線性
同餘法每一步算 $x \gets (ax + c) \bmod 2^{64}$，$a$、$c$ 是兩個選好的常數；
\texttt{unsigned long long} 的乘法和加法溢位時自動就是 $\bmod 2^{64}$，這是 C 標準規
定的行為，不是未定義行為。 低位元的週期很短，所以取最高的 $31$ 位元當鍵。

\begin{example}[label=ex:hash-probe-count]{c}{(probe\_count.c) 載入因子和平均探測次數}
#include <stdio.h>
#include <stdlib.h>
#include <math.h>

#define M 100003            // 格數（質數）
#define TRIALS 100000       // 每個載入因子做幾次失敗搜尋

// 線性同餘法的亂數產生器。種子固定，每次執行的結果都一樣。
static unsigned long long seed;

int next_key(void)
{
    seed = seed * 6364136223846793005ULL + 1442695040888963407ULL;
    return (int)(seed >> 33);           // 取最高的 31 位元
}

// 雙重雜湊的步長 h2(k) = 1 + k mod (M - 2)；線性探測的步長是 1。
int step_of(int key, int dbl)
{
    return dbl ? 1 + key % (M - 2) : 1;
}

// 找 key，回傳探測次數，*slot 是停下來的格子。-1 代表空格。
int search(const int *table, int key, int dbl, int *slot)
{
    int step = step_of(key, dbl);
    for (long long i = 0; i < M; i++) {
        *slot = (int)((key % M + i * step) % M);
        if (table[*slot] == key || table[*slot] == -1) {
            return (int)i + 1;
        }
    }
    return M;                       // 表滿了（這裡不會發生）
}

int main(void)
{
    const double alphas[] = {0.10, 0.25, 0.50, 0.75, 0.90, 0.95};
    double res[6][4];
    int *table = malloc(M * sizeof *table);
    int *keys = malloc(M * sizeof *keys);
    if (table == NULL || keys == NULL) {
        free(table);
        free(keys);
        return 1;
    }
    for (int dbl = 0; dbl <= 1; dbl++) {
        seed = 20261119ULL;             // 兩種方法用同一串鍵
        for (int j = 0; j < M; j++) {
            table[j] = -1;
        }
        int n = 0, slot;
        for (int a = 0; a < 6; a++) {
            while (n < (int)(alphas[a] * M)) {   // 插入偶數鍵
                int key = next_key() & ~1;
                search(table, key, dbl, &slot);
                if (table[slot] == -1) {   // 停在空格：不在表裡
                    table[slot] = key;
                    keys[n++] = key;
                }
            }
            long long hit = 0, miss = 0;
            for (int i = 0; i < n; i++) {    // 成功：每個鍵找一次
                hit += search(table, keys[i], dbl, &slot);
            }
            for (int t = 0; t < TRIALS; t++) {   // 失敗：奇數鍵
                miss += search(table, next_key() | 1, dbl, &slot);
            }
            res[a][2 * dbl] = (double)hit / n;
            res[a][2 * dbl + 1] = (double)miss / TRIALS;
        }
    }
    printf("        ------ linear probing ------   ");
    printf("------ double hashing ------\n");
    printf("alpha   hit   Knuth   miss   Knuth     ");
    printf("hit  uniform   miss uniform\n");
    for (int a = 0; a < 6; a++) {
        double x = alphas[a];
        printf("%.2f %6.2f %7.2f %6.2f %7.2f  ", x,
               res[a][0], 0.5 * (1 + 1 / (1 - x)),
               res[a][1], 0.5 * (1 + 1 / ((1 - x) * (1 - x))));
        printf("%6.2f %7.2f %6.2f %7.2f\n",
               res[a][2], log(1 / (1 - x)) / x,
               res[a][3], 1 / (1 - x));
    }
    free(table);
    free(keys);
    return 0;
}
\end{example}
\iosample{}{\ \ \ \ \ \ \ \ ------\ linear\ probing\ ------\ \ \ ------\ double\ hashing\ ------ \\ alpha\ \ \ hit\ \ \ Knuth\ \ \ miss\ \ \ Knuth\ \ \ \ \ hit\ \ uniform\ \ \ miss\ uniform \\ 0.10\ \ \ 1.05\ \ \ \ 1.06\ \ \ 1.12\ \ \ \ 1.12\ \ \ \ 1.05\ \ \ \ 1.05\ \ \ 1.11\ \ \ \ 1.11 \\ 0.25\ \ \ 1.17\ \ \ \ 1.17\ \ \ 1.39\ \ \ \ 1.39\ \ \ \ 1.16\ \ \ \ 1.15\ \ \ 1.33\ \ \ \ 1.33 \\ 0.50\ \ \ 1.50\ \ \ \ 1.50\ \ \ 2.49\ \ \ \ 2.50\ \ \ \ 1.41\ \ \ \ 1.39\ \ \ 1.94\ \ \ \ 2.00 \\ 0.75\ \ \ 2.50\ \ \ \ 2.50\ \ \ 8.27\ \ \ \ 8.50\ \ \ \ 1.91\ \ \ \ 1.85\ \ \ 3.80\ \ \ \ 4.00 \\ 0.90\ \ \ 5.35\ \ \ \ 5.50\ \ 48.83\ \ \ 50.50\ \ \ \ 2.66\ \ \ \ 2.56\ \ \ 9.49\ \ \ 10.00 \\ 0.95\ \ 10.00\ \ \ 10.50\ 193.24\ \ 200.50\ \ \ \ 3.28\ \ \ \ 3.15\ \ 19.29\ \ \ 20.00}

\texttt{hit} 是成功搜尋、\texttt{miss} 是失敗搜尋。 線性探測的實測值和 Knuth 的公
式相當接近，$\alpha = 0.95$ 時失敗搜尋要將近 $200$ 次，公式是 $200.5$；雙重雜湊緊跟著
均勻雜湊的公式。

失敗搜尋的實測都比公式低一些，原因不在定理~\ref{thm:hash-unsucc} 的上界：$m$ 這麼大
時，精確值 $(m+1)/(m-n+1)$ 和 $1/(1 - \alpha)$ 差不到 $0.01$。 線性探測的群聚長短起伏
很大，這裡只量了一張表；換一個種子，$\alpha = 0.95$ 的失敗搜尋從一百八十幾到兩百二十
幾都可能。 雙重雜湊偏低，是因為這裡的鍵都小於 $2^{31}$。 寫成 $k = qM + r$，
$h_1(k) = r$，而 $k \bmod (M - 2) = (2q + r) \bmod (M - 2)$ 也由同一個商 $q$ 決定，所以
$h_1$、$h_2$ 並不獨立，和均勻雜湊的假設有一點距離。 把鍵換成 $62$ 位元的亂數再量，實測
和公式就幾乎一樣。

\subsection{鏈結法和開放定址怎麼選}

鏈結法的 $\alpha$ 可以大於 $1$，變大時慢慢變慢；刪除直接從串列拿掉；代價是每個鍵多
一個指標、每次插入配置一個節點，走串列時每個節點可能在記憶體的不同地方。 開放定址只
有一個陣列，鍵小的時候省很多空間，線性探測讀的是相鄰的格子，對快取友善；代價是
$\alpha$ 接近 $1$ 時急遽變慢，刪除要用標記，標記多了要重建。 所以鍵數難預估、刪除頻
繁、每筆資料很大時用鏈結法；鍵數大致已知、資料小、以查詢為主時用開放定址。

\section{再雜湊、可延伸雜湊與通用雜湊}
\label{sec:hash-dynamic}

\subsection{再雜湊}

前面的分析都假設 $\alpha$ 是常數。 鍵一直加進來，$\alpha$ 就一直變大，鏈結法的串列
越來越長，開放定址最後會滿。 做法和第~\ref{sec:amo-dynarray}~節的動態陣列一樣：$\alpha$ 超過門檻
（例如開放定址 $0.5$、鏈結法 $1$）時，配置一張大約兩倍大的新表，把舊表的鍵逐一插入
新表，再釋放舊表。 這叫\textbf{再雜湊}（rehashing）。

搬的時候不能照格子號碼直接複製，因為 $h$ 和 $m$ 有關，$m$ 變了，每個鍵的位置都要重
算。 把第~\ref{sec:hash-open}~節線性探測的表（$m = 13$，九個鍵）搬到 $m' = 29$（大
於 $2 \times 13$ 的最小質數），依舊表的格子順序逐一插入：

\begin{center}
\begin{tabular}{lccccccccc}
\toprule
鍵 & 5347 & 6296 & 1682 & 8170 & 5906 & 3051 & 7193 & 2299 & 2014 \\
\midrule
舊位置（$m = 13$）& 4 & 5 & 6 & 7 & 8 & 9 & 10 & 11 & 12 \\
$k \bmod 29$ & 11 & 3 & 0 & 21 & 19 & 6 & 1 & 8 & 13 \\
\bottomrule
\end{tabular}
\end{center}

九個鍵在新表裡沒有任何碰撞，每個都只探測 $1$ 次，舊表裡那一長段群聚完全散開了，
$\alpha$ 降到 $9/29 \approx 0.31$。 一次再雜湊要 $\Theta(n + m)$，但和動態陣列一樣，
表變成兩倍之後要再插入大約那麼多個鍵才會再擴表，把搬動的成本平攤到這些插入上，每次
插入的平攤成本仍然是期望 $O(1)$。 再雜湊也用來清掉刪除標記（第~\ref{sec:hash-tomb}~節），
這時新表的大小可以不變。

要注意「再雜湊」這個詞在一些資料結構課本裡有另一個意思：碰撞時改用第二個、第三個
雜湊函數 $h_2, h_3, \ldots$ 算下一個位置（雙重雜湊是它的一個特例）。 考題寫「rehashing」
時，先從上下文判斷是哪一種。

\subsection{可延伸雜湊}

再雜湊一次要搬動全部的鍵。 資料存在磁碟上時，這代價太大：搬一次要把整個檔案讀寫一
遍。 \textbf{可延伸雜湊}（extendible hashing）讓表一次只長一點點，每次只拆開一個桶。

資料放在固定容量的\textbf{桶}（bucket，在磁碟上就是一頁）裡。 另外有一個
\textbf{目錄}（directory），是一個有 $2^d$ 個指標的陣列，$d$ 叫做\textbf{全域深度}
（global depth）。 每個鍵先算出一個夠長的雜湊值（下面用 $5$ 個位元），取\textbf{最
前面的 $d$ 個位元}當目錄的索引，順著指標找到桶。 幾個相鄰的目錄項可以指向同一個桶；
每個桶記一個\textbf{局部深度}（local depth）$d'$，表示這個桶裡的鍵前 $d'$ 個位元都
相同，有 $2^{d - d'}$ 個目錄項指向它。

插入時如果桶滿了，就把它拆成兩個桶，局部深度各加一，依第 $d' + 1$ 個位元把鍵分到兩
邊：

\begin{itemize}
\item $d' < d$：有不只一個目錄項指向這個桶，把其中一半改指新的桶就好，目錄不變。
\item $d' = d$：只有一個目錄項指向它，目錄先加倍（$d$ 加一，新的第 $x$ 項指向舊的第
$\lfloor x/2 \rfloor$ 項指的桶，也就是去掉最後一個位元），再照上一種情形拆。
\end{itemize}

拆完如果新的桶還是滿的（所有鍵的下一個位元都一樣），就再拆一次。 下面用容量 $2$ 的桶
示範。 七個鍵的雜湊值如下：

\begin{center}
\begin{tabular}{lccccccc}
\toprule
鍵 & A & B & C & D & E & F & G \\
\midrule
$h(k)$ & \texttt{00110} & \texttt{10011} & \texttt{01101} & \texttt{11010} &
\texttt{01001} & \texttt{11100} & \texttt{01110} \\
\bottomrule
\end{tabular}
\end{center}

放 E = \texttt{0}1001 時第一個桶滿了，而且 $d' = d = 1$：目錄先加倍成 $d = 2$
（\texttt{00}、\texttt{01} 都先指向原來的桶），再依第 $2$ 個位元拆，A = 0\texttt{0}110
留下，C、E（0\texttt{1}\ldots）到新的桶，見圖~\ref{fig:hash-extendible}(b)；B、D 的桶沒有動，局部深度仍是 $1$。
放 F = \texttt{11}100 時 B、D 的桶滿了，但 $d' = 1 < d = 2$，只拆桶、把目錄項
\texttt{11} 改指新的桶，見 (c)。 放 G = \texttt{01}110 時 C、E 的桶滿了，$d' = d = 2$，
目錄再加倍成 $d = 3$，依第 $3$ 個位元把 E = 01\texttt{0}01 和 C、G（01\texttt{1}\ldots）
分開，見 (d)；其他三個桶不動，各有兩個目錄項指著。

\begin{figure}[htbp]
\centering
\begin{tikzpicture}[x=1cm, y=0.52cm]
% (a) d = 1，放完 A、B、C、D
\begin{scope}
\node[font=\small] at (1.2, 1.1) {(a) $d = 1$};
\node[hdir] (a0) at (0, 0) {0};
\node[hdir] (a1) at (0, -1) {1};
\node[hpage] (ap) at (2.4, 0) {A C \quad $d'{=}1$};
\node[hpage] (aq) at (2.4, -1) {B D \quad $d'{=}1$};
\draw[diagarrow] (a0) -- (ap);
\draw[diagarrow] (a1) -- (aq);
\end{scope}
% (b) 放 E：目錄加倍
\begin{scope}[xshift=6.6cm]
\node[font=\small] at (1.2, 1.1) {(b) 放 E，目錄加倍，$d = 2$};
\foreach \b [count=\i from 0] in {00,01,10,11} { \node[hdir] (b\i) at (0, -\i) {\b}; }
\node[hpage] (bp) at (2.4, 0) {A \quad $d'{=}2$};
\node[hpage] (bq) at (2.4, -1) {C E \quad $d'{=}2$};
\node[hpage] (br) at (2.4, -2.5) {B D \quad $d'{=}1$};
\draw[diagarrow] (b0) -- (bp);
\draw[diagarrow] (b1) -- (bq);
\draw[diagarrow] (b2) -- (br);
\draw[diagarrow] (b3) -- (br);
\end{scope}
% (c) 放 F：拆桶，目錄不變
\begin{scope}[yshift=-2.9cm]
\node[font=\small] at (1.2, 1.1) {(c) 放 F，只拆桶，$d = 2$};
\foreach \b [count=\i from 0] in {00,01,10,11} { \node[hdir] (c\i) at (0, -\i) {\b}; }
\node[hpage] (cp) at (2.4, 0) {A \quad $d'{=}2$};
\node[hpage] (cq) at (2.4, -1) {C E \quad $d'{=}2$};
\node[hpage] (cr) at (2.4, -2) {B \quad $d'{=}2$};
\node[hpage] (cs) at (2.4, -3) {D F \quad $d'{=}2$};
\foreach \i/\t in {0/cp, 1/cq, 2/cr, 3/cs} { \draw[diagarrow] (c\i) -- (\t); }
\end{scope}
% (d) 放 G：目錄再加倍
\begin{scope}[xshift=6.6cm, yshift=-2.9cm]
\node[font=\small] at (1.2, 1.1) {(d) 放 G，目錄加倍，$d = 3$};
\foreach \b [count=\i from 0] in {000,001,010,011,100,101,110,111} {
    \node[hdir] (d\i) at (0, -\i) {\b};
}
\node[hpage] (dp) at (2.6, -0.5) {A \quad $d'{=}2$};
\node[hpage] (dq) at (2.6, -2) {E \quad $d'{=}3$};
\node[hpage] (dr) at (2.6, -3) {C G \quad $d'{=}3$};
\node[hpage] (ds) at (2.6, -4.5) {B \quad $d'{=}2$};
\node[hpage] (dt) at (2.6, -6.5) {D F \quad $d'{=}2$};
\foreach \i/\t in {0/dp, 1/dp, 2/dq, 3/dr, 4/ds, 5/ds, 6/dt, 7/dt} {
    \draw[diagarrow] (d\i.east) -- (\t.west);
}
\end{scope}
\end{tikzpicture}
\diagcaption{可延伸雜湊，桶的容量是 $2$，用雜湊值最前面的 $d$ 個位元查目錄。}
\label{fig:hash-extendible}
\end{figure}

每次只拆一個桶、只搬那個桶裡的鍵；加倍的只是目錄，目錄只存指標，常常整個放得進記憶
體，搜尋一個鍵只要讀一個桶。 代價是只要有一個桶一直被拆，目錄就得一直加倍。

\begin{examnote}
可延伸雜湊的題目，先確認用的是雜湊值的最前面幾個位元還是最後面幾個位元（本章用最前
面，也有很多課本用最後面），以及桶的容量。 每一步畫出目錄（每一項的位元字串）和每個
桶的局部深度、內容；目錄加倍和只拆桶要分清楚，判斷的依據是 $d'$ 和 $d$ 相不相等。
\end{examnote}

\subsection{通用雜湊}

任何一個固定的 $h$ 都有一批鍵全部落在同一格（第~\ref{sec:hash-func}~節）。 有人知
道伺服器用什麼雜湊函數，就能故意送一堆碰撞的鍵讓它變慢，簡單均勻雜湊的假設不成立。 \textbf{通用雜湊}（universal hashing）的想法是：先準備
一整族函數，程式開始執行時才從裡面隨機挑一個。 輸入是事先決定的，不知道會挑到哪一
個，所以沒有一組輸入對大部分的函數都很糟。 這時的期望值是對「挑函數」這件事取的，
不需要對輸入做任何假設。

\begin{definition}[通用雜湊族]
\label{def:hash-universal}
一族函數 $\mathcal{H}$，每個都把 $U$ 對應到 $\{0, \ldots, m-1\}$。 如果對任兩個不
同的鍵 $x \ne y$，滿足 $h(x) = h(y)$ 的函數 $h$ 不超過 $|\mathcal{H}|/m$ 個，就說
$\mathcal{H}$ 是\textbf{通用的}（universal）。 換句話說，從 $\mathcal{H}$ 隨機挑一
個 $h$，任兩個固定的不同鍵碰撞的機率不超過 $1/m$。
\end{definition}

\begin{theorem}
\label{thm:hash-universal}
用鏈結法，$h$ 從通用雜湊族裡隨機挑選。 對任何一組事先固定的 $n$ 個鍵，以及任何一個
鍵 $k$，第 $h(k)$ 條串列裡除了 $k$ 本身以外的鍵，期望個數不超過 $\alpha$。 所以搜尋
$k$ 期望比較的次數，$k$ 不在表裡時不超過 $\alpha$，在表裡時不超過 $1 + \alpha$。
\end{theorem}

\begin{proof}
對表裡每個不等於 $k$ 的鍵 $x$，令 $I_x = I\{h(x) = h(k)\}$。 由定義
$\mathrm{E}[I_x] = \Pr[h(x) = h(k)] \le 1/m$。 這樣的 $x$ 最多 $n$ 個，加起來期望值不
超過 $n/m = \alpha$。 搜尋 $k$ 比較的就是這條串列上的鍵；$k$ 在表裡時再加上它自己。
\end{proof}

和定理~\ref{thm:hash-chain-len} 比，結論的形式一樣，前提卻不同：這裡沒有假設鍵的分
布，隨機性完全來自程式自己挑函數。 一個常用的通用雜湊族是：取一個比所有鍵都大的質
數 $p$，
\[
h_{a,b}(k) = \bigl( (ak + b) \bmod p \bigr) \bmod m, \qquad
1 \le a \le p - 1,\ 0 \le b \le p - 1.
\]
每組 $(a, b)$ 是一個函數，一共 $p(p-1)$ 個。 它是通用的（Carter 與 Wegman，1979；
CLRS 第 4 版第 11.3.4 節有完整證明，這裡不證）。 例如 $p = 10007$、$m = 10$，執行時隨機挑到
$a = 37$、$b = 58$，則 $h(2026) = (75020 \bmod 10007) \bmod 10 = 4971 \bmod 10 = 1$。
下次執行挑到別的 $(a, b)$，同一個鍵就落在別的格子。

\section{布隆過濾器}
\label{sec:hash-bloom}

刷卡機還有一件事要做：擋下停權的卡（掛失、畢業、被盜刷）。 停權名單在學校的伺服器
上，有二十萬張；每刷一次卡就連線去問，中午排隊的人會等不及。 把名單整份放進刷卡機
也不行：用前面的雜湊表存二十萬個九位數卡號，光是鍵就要 $1.6$ MB，刷卡機只剩幾百 KB
的記憶體。 不過絕大多數刷卡的人都不在名單上，刷卡機其實只需要一個很小的東西，能很
快地說出「這張一定沒有停權」；少數說不準的，再去問伺服器。

\textbf{布隆過濾器}（Bloom filter，Burton Bloom 在 1970 年提出）做的就是這件事。 它只
回答「$x$ 在不在集合 $S$ 裡」，不存鍵本身，所以很省空間；代價是偶爾會把不在 $S$ 裡
的 $x$ 說成「可能在」。 反過來的錯誤不會發生：它說「不在」，就一定不在。

\subsection{結構與操作}

布隆過濾器是一個有 $m$ 個位元的陣列 $B[0..m-1]$，一開始全部是 $0$，再加上 $k$ 個雜湊
函數 $h_1, h_2, \ldots, h_k$，每個都把鍵對應到 $\{0, 1, \ldots, m - 1\}$。
\begin{itemize}
\item \textbf{插入} $x$：把 $B[h_1(x)], B[h_2(x)], \ldots, B[h_k(x)]$ 這 $k$ 個位元都設成
$1$（原本已經是 $1$ 的就不變）。
\item \textbf{查詢} $x$：看同樣的 $k$ 個位元。 只要有一個是 $0$，回答「不在」；全部是
$1$，回答「可能在」。
\end{itemize}

\begin{pseudo}[label=alg:hash-bloom]{布隆過濾器的插入與查詢}
\Procedure{Bloom-Insert}{$B, x$}
    \For{$i \gets 1$ \To $k$}
        \State $B[h_i(x)] \gets 1$
    \EndFor
\EndProcedure
\Function{Bloom-Query}{$B, x$}
    \For{$i \gets 1$ \To $k$}
        \If{$B[h_i(x)] = 0$}
            \State \Return \FALSE \Comment{一定不在}
        \EndIf
    \EndFor
    \State \Return \TRUE \Comment{可能在}
\EndFunction
\end{pseudo}

兩個操作都只算 $k$ 個雜湊值、看 $k$ 個位元，時間是 $O(k)$，和集合裡有幾個鍵無關。 空
間是 $m$ 個位元，也和鍵有多長無關：九位數的卡號和一百個字元的網址，佔的空間一樣。

\begin{theorem}
\label{thm:hash-bloom-noneg}
布隆過濾器不會有\textbf{偽陰性}（false negative）：如果 $x$ 已經插入過，
$\proc{Bloom-Query}(x)$ 一定回傳 \TRUE。
\end{theorem}

\begin{proof}
插入 $x$ 時，第 $3$ 行把 $B[h_1(x)], \ldots, B[h_k(x)]$ 都設成 $1$。 之後的操作只有插
入和查詢，插入只會把位元設成 $1$，查詢不改任何位元，所以這 $k$ 個位元一直是 $1$。 查
詢 $x$ 時第 $6$ 行的條件永遠不成立，迴圈跑完，回傳 \TRUE。
\end{proof}

會出錯的是另一個方向。 $x$ 沒有插入過，但它的 $k$ 個位元剛好都被別的鍵設成了 $1$，查
詢就回答「可能在」，這叫\textbf{偽陽性}（false positive）。 所以「可能在」只是可能：刷
卡機遇到「可能在」，再連線問伺服器確認；遇到「不在」就直接放行。 只要偽陽性很少，大部
分的刷卡都不必連線。

\subsection{手算例子}
\label{sec:hash-bloom-trace}

取 $m = 13$、$k = 3$。 三個雜湊函數借用第~\ref{sec:hash-open}~節雙重雜湊的式子：
\[
h_i(x) = \bigl( (x \bmod 13) + (i - 1)\,(1 + x \bmod 11) \bigr) \bmod 13, \qquad i = 1, 2, 3,
\]
也就是從 $x \bmod 13$ 出發，每次跳 $1 + (x \bmod 11)$ 格。 實務上常這樣用兩個雜湊值組出
$k$ 個，不必真的寫 $k$ 個不同的函數。 依序插入四張停權卡 $2037$、$5512$、$6180$、
$7423$，再查詢三張卡，見圖~\ref{fig:hash-bloom}。

\begin{figure}[htbp]
\centering
\begin{tikzpicture}
\foreach \j in {0,...,12} { \node[idxlbl] at ({0.6+\j*0.8}, 0.5) {\j}; }
\node[anchor=east, font=\footnotesize] at (0.2,0.00) {開始};
\node[hbit] at (0.60,0.00) {0};
\node[hbit] at (1.40,0.00) {0};
\node[hbit] at (2.20,0.00) {0};
\node[hbit] at (3.00,0.00) {0};
\node[hbit] at (3.80,0.00) {0};
\node[hbit] at (4.60,0.00) {0};
\node[hbit] at (5.40,0.00) {0};
\node[hbit] at (6.20,0.00) {0};
\node[hbit] at (7.00,0.00) {0};
\node[hbit] at (7.80,0.00) {0};
\node[hbit] at (8.60,0.00) {0};
\node[hbit] at (9.40,0.00) {0};
\node[hbit] at (10.20,0.00) {0};
\node[anchor=east, font=\footnotesize] at (0.2,-0.70) {插入 2037：9, 12, 2};
\node[hbit] at (0.60,-0.70) {0};
\node[hbit] at (1.40,-0.70) {0};
\node[hbit, hlnew] at (2.20,-0.70) {1};
\node[hbit] at (3.00,-0.70) {0};
\node[hbit] at (3.80,-0.70) {0};
\node[hbit] at (4.60,-0.70) {0};
\node[hbit] at (5.40,-0.70) {0};
\node[hbit] at (6.20,-0.70) {0};
\node[hbit] at (7.00,-0.70) {0};
\node[hbit, hlnew] at (7.80,-0.70) {1};
\node[hbit] at (8.60,-0.70) {0};
\node[hbit] at (9.40,-0.70) {0};
\node[hbit, hlnew] at (10.20,-0.70) {1};
\node[anchor=east, font=\footnotesize] at (0.2,-1.40) {插入 5512：0, 2, 4};
\node[hbit, hlnew] at (0.60,-1.40) {1};
\node[hbit] at (1.40,-1.40) {0};
\node[hbit, hlprobe] at (2.20,-1.40) {1};
\node[hbit] at (3.00,-1.40) {0};
\node[hbit, hlnew] at (3.80,-1.40) {1};
\node[hbit] at (4.60,-1.40) {0};
\node[hbit] at (5.40,-1.40) {0};
\node[hbit] at (6.20,-1.40) {0};
\node[hbit] at (7.00,-1.40) {0};
\node[hbit] at (7.80,-1.40) {1};
\node[hbit] at (8.60,-1.40) {0};
\node[hbit] at (9.40,-1.40) {0};
\node[hbit] at (10.20,-1.40) {1};
\node[anchor=east, font=\footnotesize] at (0.2,-2.10) {插入 6180：5, 2, 12};
\node[hbit] at (0.60,-2.10) {1};
\node[hbit] at (1.40,-2.10) {0};
\node[hbit, hlprobe] at (2.20,-2.10) {1};
\node[hbit] at (3.00,-2.10) {0};
\node[hbit] at (3.80,-2.10) {1};
\node[hbit, hlnew] at (4.60,-2.10) {1};
\node[hbit] at (5.40,-2.10) {0};
\node[hbit] at (6.20,-2.10) {0};
\node[hbit] at (7.00,-2.10) {0};
\node[hbit] at (7.80,-2.10) {1};
\node[hbit] at (8.60,-2.10) {0};
\node[hbit] at (9.40,-2.10) {0};
\node[hbit, hlprobe] at (10.20,-2.10) {1};
\node[anchor=east, font=\footnotesize] at (0.2,-2.80) {插入 7423：0, 10, 7};
\node[hbit, hlprobe] at (0.60,-2.80) {1};
\node[hbit] at (1.40,-2.80) {0};
\node[hbit] at (2.20,-2.80) {1};
\node[hbit] at (3.00,-2.80) {0};
\node[hbit] at (3.80,-2.80) {1};
\node[hbit] at (4.60,-2.80) {1};
\node[hbit] at (5.40,-2.80) {0};
\node[hbit, hlnew] at (6.20,-2.80) {1};
\node[hbit] at (7.00,-2.80) {0};
\node[hbit] at (7.80,-2.80) {1};
\node[hbit, hlnew] at (8.60,-2.80) {1};
\node[hbit] at (9.40,-2.80) {0};
\node[hbit] at (10.20,-2.80) {1};
\node[anchor=east, font=\footnotesize] at (0.2,-3.75) {查詢 6180：5, 2, 12};
\node[hbit] at (0.60,-3.75) {1};
\node[hbit] at (1.40,-3.75) {0};
\node[hbit, hlprobe] at (2.20,-3.75) {1};
\node[hbit] at (3.00,-3.75) {0};
\node[hbit] at (3.80,-3.75) {1};
\node[hbit, hlprobe] at (4.60,-3.75) {1};
\node[hbit] at (5.40,-3.75) {0};
\node[hbit] at (6.20,-3.75) {1};
\node[hbit] at (7.00,-3.75) {0};
\node[hbit] at (7.80,-3.75) {1};
\node[hbit] at (8.60,-3.75) {1};
\node[hbit] at (9.40,-3.75) {0};
\node[hbit, hlprobe] at (10.20,-3.75) {1};
\node[anchor=east, font=\footnotesize] at (0.2,-4.45) {查詢 3141：8, 2, 9};
\node[hbit] at (0.60,-4.45) {1};
\node[hbit] at (1.40,-4.45) {0};
\node[hbit] at (2.20,-4.45) {1};
\node[hbit] at (3.00,-4.45) {0};
\node[hbit] at (3.80,-4.45) {1};
\node[hbit] at (4.60,-4.45) {1};
\node[hbit] at (5.40,-4.45) {0};
\node[hbit] at (6.20,-4.45) {1};
\node[hbit, hbzero] at (7.00,-4.45) {0};
\node[hbit] at (7.80,-4.45) {1};
\node[hbit] at (8.60,-4.45) {1};
\node[hbit] at (9.40,-4.45) {0};
\node[hbit] at (10.20,-4.45) {1};
\node[anchor=east, font=\footnotesize] at (0.2,-5.15) {查詢 4126：5, 7, 9};
\node[hbit] at (0.60,-5.15) {1};
\node[hbit] at (1.40,-5.15) {0};
\node[hbit] at (2.20,-5.15) {1};
\node[hbit] at (3.00,-5.15) {0};
\node[hbit] at (3.80,-5.15) {1};
\node[hbit, hlprobe] at (4.60,-5.15) {1};
\node[hbit] at (5.40,-5.15) {0};
\node[hbit, hlprobe] at (6.20,-5.15) {1};
\node[hbit] at (7.00,-5.15) {0};
\node[hbit, hlprobe] at (7.80,-5.15) {1};
\node[hbit] at (8.60,-5.15) {1};
\node[hbit] at (9.40,-5.15) {0};
\node[hbit] at (10.20,-5.15) {1};

\end{tikzpicture}
\diagcaption{布隆過濾器（$m = 13$，$k = 3$）。 左邊是操作和算出的三個位置。 插入時橘框
是這次由 $0$ 變 $1$ 的位元，粉紅底是本來就是 $1$ 的；查詢時粉紅底是看過、是 $1$ 的位元，
深色底是看到的 $0$，看到它就停。}
\label{fig:hash-bloom}
\end{figure}

以 $2037$ 為例：$2037 = 13 \times 156 + 9$，$2037 = 11 \times 185 + 2$，所以從第 $9$ 格出
發、每次跳 $3$ 格，三個位置是 $9$、$12$、$15 \bmod 13 = 2$。 其他三張卡依同樣的算法得到
$\{0, 2, 4\}$、$\{5, 2, 12\}$、$\{0, 10, 7\}$。 四張卡一共設了 $12$ 次位元，但第 $0$、
$2$、$12$ 格被設了不只一次，最後只有 $8$ 個位元是 $1$。

三個查詢：
\begin{itemize}
\item $6180$ 插入過，三個位元 $5, 2, 12$ 都是 $1$，回答「可能在」。 這是對的。
\item $3141$：$3141 \bmod 13 = 8$，$1 + 3141 \bmod 11 = 8$，位置是 $8, 2, 9$。 第 $8$ 格
是 $0$，回答「不在」。 看到第一個 $0$ 就可以停，不必看完三格。
\item $4126$：$4126 \bmod 13 = 5$，$1 + 4126 \bmod 11 = 2$，位置是 $5, 7, 9$，分別被
$6180$、$7423$、$2037$ 設成 $1$，回答「可能在」。 但 $4126$ 從來沒有插入過，這是一次偽
陽性。
\end{itemize}

這個例子的過濾器已經有 $8/13$ 的位元是 $1$，太滿了。 把 $1000$ 到 $9999$ 之間沒插入過
的卡號都查一次，大約有五分之一會被誤判成「可能在」。 下一小節會算出這個比例怎麼由 $m$、
$n$、$k$ 決定。

\subsection{偽陽性率}
\label{sec:hash-bloom-fp}

和第~\ref{sec:hash-chain-analysis}~節一樣，先假設雜湊值像亂數：每個 $h_i(x)$ 都等可能
是 $0$ 到 $m - 1$ 的任何一格，而且不同的 $i$、不同的鍵之間互相獨立。 插入了 $n$ 個鍵之
後，查詢一個不在集合裡的鍵 $y$，回答「可能在」的機率是多少？

先看一個固定的位元 $B[j]$。 $n$ 個鍵一共算了 $kn$ 個雜湊值，每一個落在 $j$ 以外的機率
是 $1 - 1/m$，而且互相獨立，所以
\[
\Pr[B[j] = 0] = \left(1 - \frac{1}{m}\right)^{kn} \approx e^{-kn/m}.
\]
近似的根據是引理~\ref{lem:mth-exp}：把 $(1 - 1/m)^{kn}$ 寫成
$\bigl((1 - 1/m)^m\bigr)^{kn/m}$，括號裡的數夾在 $e^{-1}(1 - 1/m)$ 和 $e^{-1}$ 之間，$m$
很大時就是 $e^{-1}$。 記 $p = e^{-kn/m}$，表示大約有 $p$ 比例的位元還是 $0$。

$y$ 的 $k$ 個位置也是均勻亂選的，每一個看到 $1$ 的機率大約是 $1 - p$。 $k$ 個都看到 $1$
才會誤判，把它們當成互相獨立，得到偽陽性率
\begin{equation}
\label{eq:hash-bloom-fp}
\varepsilon \approx (1 - p)^k = \bigl(1 - e^{-kn/m}\bigr)^k.
\end{equation}
嚴格說，各個位元是不是 $0$ 並不完全獨立（一個位元是 $1$，別的位元是 $1$ 的機會就小一
點），所以這是近似值；$m$ 很大時，$1$ 的比例會非常接近它的期望值，這個近似相當準。
第~\ref{sec:hash-bloom-c}~節的程式會拿實測值來比。

用手算的例子驗證：$m = 13$、$n = 4$、$k = 3$，式~\eqref{eq:hash-bloom-fp} 給
$(1 - e^{-12/13})^3 \approx 0.22$，和上一小節「大約五分之一」相近。

\paragraph{$k$ 要取多少？}
$m$、$n$ 固定時，$k$ 太小，每次只檢查少數幾個位元，容易全部碰巧是 $1$；$k$ 太大，每個
鍵設很多位元，陣列很快就被 $1$ 塞滿。 中間有一個最好的 $k$。

\begin{theorem}
\label{thm:hash-bloom-k}
把 $k$ 當成正實數，$\bigl(1 - e^{-kn/m}\bigr)^k$ 在
\[
k^{*} = \frac{m}{n} \ln 2 \approx 0.693\,\frac{m}{n}
\]
時最小。 這時 $p = e^{-k^{*} n/m} = \frac12$，也就是大約一半的位元是 $0$，偽陽性率是
\[
\left(\frac12\right)^{k^{*}} = 2^{-(m/n)\ln 2} \approx 0.6185^{\,m/n}.
\]
\end{theorem}

\begin{whybox}[定理~\ref{thm:hash-bloom-k} 的證明]
取對數比較好處理：$\ln \varepsilon = k \ln(1 - e^{-kn/m})$。 改用 $p = e^{-kn/m}$ 當變
數，$k$ 從 $0$ 變到 $\infty$ 時 $p$ 從 $1$ 遞減到 $0$，而且 $k = -\frac{m}{n} \ln p$，代進去
\[
\ln \varepsilon = -\frac{m}{n}\,\ln p \cdot \ln(1 - p).
\]
$m/n$ 是固定的正數，所以要讓 $\varepsilon$ 最小，就是要讓 $\varphi(p) = \ln p \cdot \ln(1 - p)$
在 $0 < p < 1$ 上最大（兩個對數都是負的，$\varphi > 0$）。 微分：
\[
\varphi'(p) = \frac{\ln(1 - p)}{p} - \frac{\ln p}{1 - p}
= \frac{\chi(p)}{p(1 - p)}, \qquad
\chi(p) = (1 - p)\ln(1 - p) - p \ln p .
\]
分母是正的，只要看 $\chi$ 的正負。 $\chi(\frac12) = 0$，$\chi(1 - p) = -\chi(p)$，而且
$p \to 0^{+}$ 時 $p \ln p \to 0$，所以 $\chi(p) \to 0$。 再微分兩次：
\[
\chi'(p) = -\ln(1 - p) - 1 - \ln p - 1, \qquad
\chi''(p) = -\frac{1}{p} + \frac{1}{1 - p} = \frac{2p - 1}{p(1 - p)}.
\]
$0 < p < \frac12$ 時 $\chi'' < 0$，$\chi$ 是凹函數，圖形在兩端點 $(0, 0)$、
$(\frac12, 0)$ 的連線之上，所以 $\chi(p) > 0$；由 $\chi(1 - p) = -\chi(p)$，
$\frac12 < p < 1$ 時 $\chi(p) < 0$。 於是 $\varphi$ 在 $p < \frac12$ 遞增、$p > \frac12$
遞減，最大值在 $p = \frac12$。 $e^{-kn/m} = \frac12$ 解出 $k = \frac{m}{n}\ln 2$；代回
去，$\varepsilon = (1 - \frac12)^{k^{*}} = 2^{-k^{*}}$，而 $2^{-\ln 2} = e^{-(\ln 2)^2}
\approx 0.6185$。
\end{whybox}

$k$ 必須是整數。 由證明，$\varepsilon$ 在 $k^{*}$ 左邊遞減、右邊遞增，所以只要比較
$\lfloor k^{*} \rfloor$ 和 $\lceil k^{*} \rceil$ 兩個值，取比較小的那個。 下表是每個鍵分
到 $m/n$ 個位元時的結果。

\begin{center}
\begin{tabular}{cccc}
\toprule
每個鍵的位元數 $m/n$ & $k^{*} = \frac{m}{n}\ln 2$ & 最好的整數 $k$ & 偽陽性率 \\
\midrule
$4$ & $2.77$ & $3$ & $14.7\%$ \\
$8$ & $5.55$ & $6$ & $2.16\%$ \\
$10$ & $6.93$ & $7$ & $0.82\%$ \\
$16$ & $11.09$ & $11$ & $0.046\%$ \\
\bottomrule
\end{tabular}
\end{center}

反過來，想要偽陽性率是 $\varepsilon$，令 $2^{-k^{*}} = \varepsilon$，得到
$k^{*} = \lg(1/\varepsilon)$，每個鍵需要
\[
\frac{m}{n} = \frac{k^{*}}{\ln 2} = \frac{\lg(1/\varepsilon)}{\ln 2} \approx 1.44 \lg \frac{1}{\varepsilon}
\]
個位元。 $\varepsilon = 1\%$ 時大約是 $9.6$ 個位元、$k = 7$。 回到章首的刷卡機：二十萬張
停權卡、每張 $10$ 個位元，一共兩百萬個位元，大約 $244$ KB，放得進刷卡機；沒有停權的
卡，每刷一百次平均不到一次要連線確認。 注意這個數字和鍵有多長完全無關。

\subsection{不能刪除}

卡片解除停權時，能不能把它的 $k$ 個位元清成 $0$？ 不行。 那些位元可能也屬於別的鍵。
在圖~\ref{fig:hash-bloom} 的過濾器裡「刪除」$6180$，把第 $5$、$2$、$12$ 格清成 $0$；
再查 $2037$，它的位置是 $9, 12, 2$，第 $12$ 格已經是 $0$，回答「不在」。 $2037$ 明明
還在名單上，這是一次偽陰性，定理~\ref{thm:hash-bloom-noneg} 的保證就沒了。 這和開放
定址不能把刪掉的格子改回空的（第~\ref{sec:hash-tomb}~節）是同一類問題：一個位置被好幾
個鍵共用，清掉它會影響別人。

需要刪除時可以用\textbf{計數型布隆過濾器}（counting Bloom filter）：每個位置從一個位元
改成一個小計數器，插入時 $k$ 個計數器各加一，刪除時各減一，查詢時看 $k$ 個計數器是不是
都大於 $0$，代價是空間變成好幾倍。 名單變動不頻繁時，另一個做法是定期用最新的名單重建
整個過濾器，這和第~\ref{sec:hash-tomb}~節清理刪除標記的想法一樣。

\begin{examnote}
布隆過濾器的題目常見三種：(1) 給位元陣列和雜湊函數，判斷插入會不會改到位元、查詢回答
什麼，逐一算出 $k$ 個位置再看位元，有一個 $0$ 就是「不在」；(2) 問最佳的雜湊函數個數，
答 $\frac{m}{n}\ln 2$，題目的選項常寫成「$3\ln 2$」這種形式，先認清 $m$（位元數）和 $n$
（鍵數）；(3) 問偽陽性率，寫出 $\bigl(1 - (1 - \frac1m)^{kn}\bigr)^k \approx
\bigl(1 - e^{-kn/m}\bigr)^k$ 兩種寫法都可以。 題目說「刪除某個鍵就把它的位元清成 $0$」，
幾乎都是要你指出這樣會造成偽陰性。
\end{examnote}

\subsection{C 實作與實測}
\label{sec:hash-bloom-c}

下面的程式先重做圖~\ref{fig:hash-bloom} 的例子，再做一個實驗：$m = 10007$ 個位元、
$n = 1000$ 個鍵（每個鍵大約 $10$ 個位元），$k$ 從 $1$ 試到 $10$，每個 $k$ 都查詢
$100000$ 個不在集合裡的鍵，算出誤判的比例，旁邊印出式~\eqref{eq:hash-bloom-fp} 的值。
雜湊函數用和手算例子相同的寫法，$m = 13$ 時 $m - 2 = 11$，正好是上面的式子。 鍵用
第~\ref{sec:hash-measure}~節的線性同餘亂數產生，插入的都是偶數、查詢的都是奇數，所以查
詢的鍵一定不在集合裡。

為了好讀，程式用一個 \texttt{unsigned char} 存一個位元，浪費了八倍的空間。 真正的實作
把八個位元塞進一個位元組，用《C 程式設計基礎》第~12~章的位元運算讀寫：第 $j$ 個位元
在 \texttt{bit[j / 8]} 的第 \texttt{j \% 8} 位。

\begin{example}[label=ex:hash-bloom]{c}{(bloom\_filter.c) 布隆過濾器：手算例子與偽陽性率實測}
#include <stdio.h>
#include <stdlib.h>
#include <math.h>

// Bloom filter：m 個位元，k 個雜湊函數。
// 第 i 個函數（i 從 0 開始）用雙重雜湊的式子：
//   g_i(x) = (x mod m + i * (1 + x mod (m - 2))) mod m
typedef struct {
    unsigned char *bit;     // 一個位元用一個 unsigned char 存
    int m, k;
} Bloom;

int bloom_init(Bloom *b, int m, int k)
{
    b->bit = calloc(m, 1);              // 全部是 0
    b->m = m;
    b->k = k;
    return b->bit != NULL;
}

int g(const Bloom *b, long long x, int i)
{
    long long step = 1 + x % (b->m - 2);
    return (int)((x % b->m + i * step) % b->m);
}

void bloom_insert(Bloom *b, long long x)
{
    for (int i = 0; i < b->k; i++) {
        b->bit[g(b, x, i)] = 1;
    }
}

// 回傳 1 表示「可能在」，0 表示「一定不在」。
int bloom_query(const Bloom *b, long long x)
{
    for (int i = 0; i < b->k; i++) {
        if (b->bit[g(b, x, i)] == 0) {
            return 0;
        }
    }
    return 1;
}

// 和 probe_count.c 相同的線性同餘亂數，種子固定。
static unsigned long long seed;

long long next_key(void)
{
    seed = seed * 6364136223846793005ULL + 1442695040888963407ULL;
    return (long long)(seed >> 33);
}

int main(void)
{
    // 手算的例子：m = 13，k = 3，四個停權卡號。
    const long long banned[] = {2037, 5512, 6180, 7423};
    const long long ask[] = {6180, 3141, 4126};
    Bloom b;
    if (!bloom_init(&b, 13, 3)) {
        return 1;
    }
    for (int i = 0; i < 4; i++) {
        bloom_insert(&b, banned[i]);
    }
    printf("bits:");
    for (int j = 0; j < b.m; j++) {
        printf(" %d", b.bit[j]);
    }
    printf("\n");
    for (int i = 0; i < 3; i++) {
        printf("%lld:", ask[i]);
        for (int t = 0; t < b.k; t++) {
            printf(" %d", g(&b, ask[i], t));
        }
        printf(" -> %s\n", bloom_query(&b, ask[i]) ? "maybe" : "no");
    }
    free(b.bit);

    // 實測：m = 10007，放 n = 1000 個偶數鍵，再問 100000 個奇數鍵。
    const int m = 10007, n = 1000, trials = 100000;
    printf("m = %d, n = %d\n", m, n);
    printf(" k  measured  formula\n");
    for (int k = 1; k <= 10; k++) {
        if (!bloom_init(&b, m, k)) {
            return 1;
        }
        seed = 20261119ULL;             // 每個 k 用同一串鍵
        for (int i = 0; i < n; i++) {
            bloom_insert(&b, next_key() & ~1LL);
        }
        int fp = 0;                     // 奇數鍵一定不在集合裡
        for (int t = 0; t < trials; t++) {
            fp += bloom_query(&b, next_key() | 1);
        }
        double p = pow(1 - exp(-(double)k * n / m), k);
        printf("%2d  %8.4f  %7.4f\n", k, (double)fp / trials, p);
        free(b.bit);
    }
    return 0;
}
\end{example}
\iosample{}{bits: 1 0 1 0 1 1 0 1 0 1 1 0 1 \\ 6180: 5 2 12 -> maybe \\ 3141: 8 2 9 -> no \\ 4126: 5 7 9 -> maybe \\ m = 10007, n = 1000 \\ \ k\ \ measured\ \ formula \\ \ 1\ \ \ \ 0.0959\ \ \ 0.0951 \\ \ 2\ \ \ \ 0.0331\ \ \ 0.0328 \\ \ 3\ \ \ \ 0.0185\ \ \ 0.0174 \\ \ 4\ \ \ \ 0.0120\ \ \ 0.0118 \\ \ 5\ \ \ \ 0.0101\ \ \ 0.0094 \\ \ 6\ \ \ \ 0.0089\ \ \ 0.0084 \\ \ 7\ \ \ \ 0.0086\ \ \ 0.0082 \\ \ 8\ \ \ \ 0.0086\ \ \ 0.0084 \\ \ 9\ \ \ \ 0.0096\ \ \ 0.0091 \\ 10\ \ \ \ 0.0109\ \ \ 0.0101}

前四行就是圖~\ref{fig:hash-bloom}：位元陣列，和三個查詢各自算出的位置與回答。 實測的
部分，$m/n \approx 10$，由定理~\ref{thm:hash-bloom-k}，最好的 $k$ 是 $10.007 \ln 2
\approx 6.94$，取 $k = 7$，公式值 $0.82\%$，實測 $0.86\%$；$k = 1$ 時偽陽性率將近
$10\%$，$k = 10$ 又回升到 $1\%$ 以上，和上面「$k$ 太小、太大都不好」的分析一致。 實測都
比公式高一點，原因是公式把各個位元當成互相獨立，又假設 $k$ 個位置完全亂選；這裡的 $k$
個位置由兩個雜湊值組出來，也不是完全獨立。 差距在 $10\%$ 以內，拿來決定 $m$ 和 $k$ 已
經夠用。

最後把布隆過濾器和雜湊表放在一起比。 雜湊表存了鍵本身，能回答「在不在」而且不會錯，
能刪除，也能把鍵列出來；布隆過濾器每個鍵只用幾個位元，和鍵的長度無關，但只能回答
「一定不在」或「可能在」，不能刪除，也列不出存了哪些鍵。 適合它的場合是：集合很大、
大部分查詢的答案是「不在」、偶爾誤判只是多花一點時間（多問一次伺服器、多讀一次磁碟），
不會造成錯誤的結果。

\section{本章常見錯誤}
\label{sec:hash-pitfalls}

\begin{center}
\small
\begin{tabular}{p{0.24\linewidth}p{0.50\linewidth}p{0.16\linewidth}}
\toprule
錯誤 & 例子 & 見 \\
\midrule
除法的 $m$ 沒選好 & $m = 100$ 只用到最後兩位；鍵都是 $4$ 的倍數時用 $m = 12$ &
第~\ref{sec:hash-func}~節 \\
負數取餘數 & \texttt{-7 \% 13} 是 $-7$；要寫 \texttt{((k \% m) + m) \% m} &
第~\ref{sec:hash-func}~節 \\
探測次數少算一次 & 只數碰撞，沒算最後放進去或停下來的那一格 &
第~\ref{sec:hash-open}~節 \\
二次探測當成走得遍 & $m$ 不是質數，或 $\alpha > 1/2$，以為一定插得進去 &
第~\ref{sec:hash-open}~節 \\
$h_2(k)$ 選錯 & $h_2(k) = k \bmod 7$ 可能是 $0$，原地不動；$m = 12$、$h_2(k) = 4$
不互質，只走得到三格 & 第~\ref{sec:hash-open}~節 \\
刪除時清成空格 & 之後經過這一格才放進去的鍵都找不到了 & 第~\ref{sec:hash-tomb}~節 \\
插入遇到刪除標記就放 & 沒有確認後面有沒有同一個鍵，表裡出現兩個 $k$ &
第~\ref{sec:hash-tomb}~節 \\
公式張冠李戴 & 線性探測套用 $1/(1-\alpha)$；成功和失敗的公式對調 &
第~\ref{sec:hash-analysis}~節 \\
布隆過濾器清位元 & 刪除時把鍵的 $k$ 個位元清成 $0$；這些位元可能也屬於別的鍵，之後查那個鍵會
答「不在」 & 第~\ref{sec:hash-bloom}~節 \\
「可能在」當成「在」 & 布隆過濾器只保證「不在」一定對；回答「可能在」還要再確認 &
第~\ref{sec:hash-bloom}~節 \\
\bottomrule
\end{tabular}
\end{center}

\section{觀念題}
\label{sec:hash-exercises}

題型和研究所考試、期中期末考的筆試相近，有選擇、計算、追蹤和證明。 除非題目另外說
明，探測次數照本章的數法（看過幾格就算幾次，包括最後一格），$i$ 從 $0$ 開始。 解答
在本章最後。

\begin{enumerate}
\item （選擇，可能不只一個答案）下列哪些敘述正確？
\begin{enumerate}[nosep]
\item 鏈結法的載入因子可以大於 $1$。
\item 線性探測裡，起始位置不同的兩個鍵，探測序列不會經過同一格。
\item 二次探測沒有主要群聚，但起始位置相同的鍵，探測序列完全相同。
\item $m = 16$ 的雙重雜湊，只要 $h_2(k)$ 一定是奇數，每個鍵的探測序列都會走遍 $16$ 格。
\item 開放定址刪除一個鍵時，把那一格改回空的，之後的搜尋仍然正確。
\item 在均勻雜湊下，$\alpha = 0.8$ 的表做一次失敗搜尋，期望探測次數小於 $5$。
\end{enumerate}

\item （計算）鍵 $k = 6083$。
\begin{enumerate}[nosep]
\item 用除法，$m = 97$，求 $h(k)$。
\item 用乘法，$m = 1000$，$A = 0.6180339887$，求 $h(k)$。
\item 用中間平方法，$m = 100$：$k^2$ 補成八位數，取第 $4$、$5$ 位。
\item 另一個鍵 $526803147$ 切成三位數一段，$m = 1000$，分別用移位摺疊和邊界摺疊求
$h$。
\end{enumerate}

\item （追蹤）$m = 11$，$h(k) = k \bmod 11$，用線性探測依序插入
$63, 85, 75, 54, 96, 24, 33$。
\begin{enumerate}[nosep]
\item 寫出每個鍵試過的格子和最後的表。
\item 成功搜尋的平均探測次數是多少？ 假設起始位置是 $0$ 到 $10$ 的機會一樣，失敗搜
尋的平均探測次數是多少？
\item 接著刪除 $54$（用刪除標記），再依序搜尋 $33$、搜尋 $96$、插入 $44$。 寫出每
個操作的探測次數，以及 $44$ 放在哪一格。
\end{enumerate}

\item （追蹤）用二次探測依序插入 $5, 16, 27, 38, 49, 60, 71$，$m = 11$，
$h(k, i) = (k \bmod 11 + i^2) \bmod 11$。 寫出每個鍵試過的格子。 最後一個鍵插得進去嗎？ 這時表裡
還有幾個空格？ 用定理~\ref{thm:hash-quad} 解釋。

\item （追蹤）$m = 11$，雙重雜湊 $h_1(k) = k \bmod 11$、
$h_2(k) = 1 + (\lfloor k / 11 \rfloor \bmod 10)$，依序插入第 3 題的七個鍵。 寫出每個
鍵的 $h_2$、試過的格子和最後的表，並和第 3 題比較探測的總次數。

\item （計算）$m = 8$，$h(k, i) = (h(k) + i^2) \bmod 8$。
\begin{enumerate}[nosep]
\item 列出 $i = 0, 1, \ldots, 7$ 的 $i^2 \bmod 8$。 起始位置是 $0$ 的鍵，最多試得到
幾個不同的格子？
\item 舉一個例子：表裡還有空格，插入卻失敗。
\item 改成 $h(k, i) = (h(k) + i(i+1)/2) \bmod 8$，重做 (a)。
\end{enumerate}

\item （證明）$m = 2^p$，$p \ge 1$。 證明 $h(k, i) = (h(k) + i(i+1)/2) \bmod m$，
$i = 0, 1, \ldots, m - 1$ 兩兩不同。（提示：$j - i$ 和 $i + j + 1$ 一個是奇數、一個是
偶數。）

\item （計算）鏈結法，$m = 100$ 格、$n = 100$ 個鍵，簡單均勻雜湊。
\begin{enumerate}[nosep]
\item 空的格子期望有幾個？（提示：對每一格寫一個「這一格是空的」的指示變數。）
\item 成功搜尋和失敗搜尋的期望比較次數各是多少？
\item 平均每條串列長 $1$，為什麼還有大約三分之一的格子是空的？
\end{enumerate}

\item （計算）希望失敗搜尋的期望探測次數不超過 $3$。 用均勻雜湊的 $1/(1-\alpha)$ 和
線性探測的 Knuth 公式，各求 $\alpha$ 最大可以是多少。 $m = 1000$ 時，兩種情形各大
約可以放幾個鍵？ 和範例~\ref{ex:hash-probe-count} 的輸出比對。

\item （追蹤）接著圖~\ref{fig:hash-extendible}(d) 的狀態，依序插入
H（\texttt{00011}）和 I（\texttt{00101}）。 每插入一個，說明有沒有拆桶、目錄有沒有加
倍，並畫出最後的目錄和桶。 最後有幾個目錄項指向 D、F 那個桶？

\item （計算）$U = \{a, b, c, d\}$，$m = 2$，$\mathcal{H} = \{h_1, h_2, h_3\}$：
\begin{center}
\begin{tabular}{ccccc}
\toprule
 & $a$ & $b$ & $c$ & $d$ \\
\midrule
$h_1$ & 0 & 0 & 1 & 1 \\
$h_2$ & 0 & 1 & 0 & 1 \\
$h_3$ & 1 & 0 & 0 & 1 \\
\bottomrule
\end{tabular}
\end{center}
\begin{enumerate}[nosep]
\item $\mathcal{H}$ 是通用的嗎？
\item 把 $h_3$ 換成 $h_3'(a, b, c, d) = (0, 0, 0, 1)$，還是通用的嗎？
\end{enumerate}

\item \label{q:hash-bloom-trace}（追蹤）布隆過濾器 $m = 11$、$k = 2$，$h_1(x) = x \bmod 11$，
$h_2(x) = \lfloor x / 11 \rfloor \bmod 11$，依序插入 $25, 47, 63, 90$。
\begin{enumerate}[nosep]
\item 寫出每個鍵的兩個位置和最後的位元陣列 $B[0..10]$。
\item 查詢 $36$、$58$、$80$，各回答什麼？ 哪些是偽陽性？
\item 想刪除 $63$，把它的兩個位元清成 $0$。 之後查詢哪一個插入過的鍵會答錯？
\item 改成計數型布隆過濾器，寫出插入四個鍵之後的計數器 $C[0..10]$，以及刪除 $63$
（兩個計數器各減一）之後的計數器。 這時查詢 $90$、$58$ 各回答什麼？
\end{enumerate}

\item （計算）布隆過濾器要存 $n = 100$ 個鍵，用 $m = 1200$ 個位元。 用
式~\eqref{eq:hash-bloom-fp} 估計。
\begin{enumerate}[nosep]
\item 最好的整數 $k$ 是多少？ 這時偽陽性率大約多少？ 大約多少比例的位元是 $0$？
\item 想讓偽陽性率降到 $0.1\%$ 以下，$n = 100$ 時大約要多少個位元？ $k$ 取多少？
\item 有人說「$k$ 越大，查詢時檢查的位元越多，誤判一定越少」。 用 (a) 的數字說明這句話
哪裡不對。
\end{enumerate}

\end{enumerate}

\section{歷屆試題}
\label{sec:hash-pastexam}

以下是研究所入學考試的原題，題目照原卷（語言、記號、配分照錄）。 詳解在本章最後。

\begin{pastexam}[label={px:hash-1}]{成大資訊聯招 110 學年度 第 3 題}
(10\%) Assume that you have to store the four strings: ``abcdef'', ``bcdefa'', ``cdefab'', and ``defabc'', in the hash table by using the hashing technique.
\begin{enumerate}[label=(\arabic*),itemsep=0pt,topsep=2pt]
\item (5\%) Consider the hash function: the index for a specific string will be equal to the sum of the ASCII values of the characters modulo 11. Please determine how many collisions if we use this hash function.
\item (5\%) Consider the hash function: the index for a specific string will be equal to sum of ASCII values of characters multiplied by their respective order in the string after which it is modulo with 11. Please determine how many collisions if we use this hash function.
\end{enumerate}
Note: The ASCII values of a, b, c, d, e, and f are 97, 98, 99, 100, 101, and 102 respectively.
\end{pastexam}

\begin{pastexam}[label={px:hash-2}]{政大資科所 115 學年度 第 13 題}
Consider a hash table of size $M=20$ (slots indexed from 0 to 19). The primary hash function is defined as: $h_1(k)=k \mathit{\ mod\ } 20$. The secondary hash function is defined as: $h_2(k)=(k \mathit{\ mod\ } 7)+1$. Assume the hash table is initially empty before each part of the question. We will insert the following sequence of keys S in the order given:
\[ S=[8,38,3,5,28,18,65,83,25] \]
(a) (5\%) Assuming we use double hashing for collision resolution, determine the final structure of the hash table after all keys in S have been inserted. Note: The probe sequence is given by:
\[ h(k,i)=\big(h_1(k)+i\cdot h_2(k)\big) \mathit{\ mod\ } M \]
where $i$ is the probe number.

(b) (5\%) If we search for 48, what is the total number of probes made during the search?

(c) (10\%) Suppose we are given $n$ keys, $m$ hash table slots, and two uniform hash functions $h_a$ and $h_b$. Further suppose hashing scheme uses $h_a$ for the odd keys and $h_b$ for the even keys. What is the expected number of keys in a slot?
\end{pastexam}

\begin{pastexam}[label={px:hash-3}]{中央資工所 115 學年度 第 12 題}
（多選題，每一個選項單獨計分，每題 5 分）
Given an empty 11-bucket hash table with hash function h(k)=k\%11. Below is the resulting hash table after inserting 8 items.
\begin{center}
\begin{tabular}{|l|l|}
\hline
0 & \\ \hline
1 & \\ \hline
2 & 24 \\ \hline
3 & 57 \\ \hline
4 & 37 \\ \hline
5 & 27 \\ \hline
6 & 16 \\ \hline
7 & 17 \\ \hline
8 & 15 \\ \hline
9 & 38 \\ \hline
10 & \\ \hline
\end{tabular}
\end{center}
Which of the following statements is/are true?
\begin{enumerate}[label=\Alph*.]
\item 27 37 16 17 24 15 57 38 is a possible input sequence when quadratic probing is adopted.
\item 37 27 16 24 15 57 38 17 is a possible input sequence when quadratic probing is adopted.
\item 37 27 16 24 57 17 15 38 is a possible input sequence when liner probing is adopted.
\item 27 16 37 38 57 15 17 24 is a possible input sequence when liner probing is adopted.
\end{enumerate}
\end{pastexam}

\begin{pastexam}[label={px:hash-4}]{臺大資工所 115 學年度 第 15 題}
(4 points) Most hash table implementations (such as Java's \texttt{HashMap}) resize the underlying array when the load factor $\alpha$ exceeds a certain threshold (typically 0.75). Why is this resizing operation necessary for a hash table that uses separate chaining?
\begin{enumerate}[label=(\Alph*),leftmargin=2.2em,itemsep=1pt]
\item \textbf{Saturation:} To prevent the table from reaching a load factor of 1.0, at which point inserting a new element would be impossible.
\item \textbf{Ordering:} To re-sort the keys alphabetically during the transfer, enabling the use of binary search for future lookups.
\item \textbf{Complexity:} To prevent the average length of the linked lists (collision chains) from growing linearly with the number of elements, thereby maintaining $O(1)$ average access time.
\item \textbf{Uniqueness:} To redistribute the keys across a larger range so that all existing collisions are permanently eliminated.
\item \textbf{Memory contiguity:} To ensure that all elements in the linked lists are stored in adjacent memory addresses, minimizing CPU cache misses and memory fragmentation.
\end{enumerate}
\end{pastexam}

\begin{pastexam}[label={px:hash-5}]{成大資訊聯招 113 學年度 第 5 題}
\textbf{[3\%, 2\%]} Given the following bloom filter.
\begin{center}
\small
\begin{tabular}{|*{15}{c|}}
\hline
[0] & [1] & [2] & [3] & [4] & [5] & [6] & [7] & [8] & [9] & [10] & [11] & [12] & [13] & [14] \\ \hline
1 & 1 & 0 & 1 & 1 & 1 & 1 & 1 & 0 & 1 & 0 & 0 & 0 & 1 & 0 \\ \hline
\end{tabular}
\end{center}
\begin{enumerate}[label=\roman*)]
\item Assume that the bloom filter has 3 hash functions $f_1(k)$, $f_2(k)$, and $f_3(k)$:
\[ f_1(k)=(3k) \bmod m,\qquad f_2(k)=(2k+1) \bmod m,\qquad f_3(k)=k^2 \bmod m, \]
where $k$ represents the key and $m$ is the size of bit array for the bloom filter. Which of the following statements is true?
\begin{enumerate}[label=(\Alph*)]
\item When inserting $k$=12, the bits in the bloom filter are not changed.
\item To delete $k$=2, the bits at indices 4, 5, and 6 are changed to 0.
\item When querying the bloom filter for $k$=15, the answer is ``yes.''
\item When querying the bloom filter for $k$=8, the answer is ``yes.''
\item None of the above.
\end{enumerate}
\item Considering that the elements in the bloom filter is 5. To minimize the false positive probability, what is the ideal optimum number of hash functions?
\begin{enumerate}[label=(\Alph*)]
\item 15
\item 5
\item 3*ln(2)
\item 5*ln(2)
\item None of the above
\end{enumerate}
\end{enumerate}
\end{pastexam}


\section{實作練習}
\label{sec:hash-practice}

下面是 Aki-Judge 上和本章相關的題目。

\begin{center}
\begin{tabular}{llc}
\toprule
小節 & 題目 & 難度 \\
\midrule
鏈結法 & \akiref{integer-hash-set-with-chaining}{用鏈結法實作整數集合} & 基本 \\
鏈結法 & \akiref{integer-hash-map-with-chaining}{用鏈結法實作整數對映} & 基本 \\
開放定址 & \akiref{open-addressing-integer-set}{線性探測的整數集合} & 基本 \\
刪除標記 & \akiref{linear-probing-with-deletion-marks}{線性探測雜湊表與刪除標記} & 進階 \\
再雜湊 & \akiref{parking-plate-registry-with-move-to-front-chains}{停車場車牌登記：前移串鏈與擴表} & 進階 \\
應用 & \akiref{number-of-distinct-values}{不同數值的個數} & 暖身 \\
應用 & \akiref{most-frequent-value}{出現次數最多的數} & 暖身 \\
應用 & \akiref{first-value-appearing-once}{第一個只出現一次的數} & 暖身 \\
\bottomrule
\end{tabular}
\end{center}

前三題只交函式，位置都規定用
\texttt{((key \% M) + M) \% M} 算，鍵可能是負數。 「線性探測雜湊表與刪除標記」的
規則和第~\ref{sec:hash-tomb}~節相同，可以先用圖~\ref{fig:hash-tomb} 的例子對一次。 最後三題是計數：把範例~\ref{ex:hash-chain} 的「餘額」改成「出現次數」就能做。

\section{讀完你應該會}
\label{sec:hash-checklist}

\begin{itemize}
\item 能用除法、乘法、中間平方、摺疊算出雜湊值，並說明除法的 $m$ 為什麼要選質數。
\item 能寫出鏈結法的三個操作，並用指示變數推導成功與失敗搜尋的期望比較次數。
\item 能對一串鍵做線性探測、二次探測、雙重雜湊的追蹤，寫出每個鍵試過的格子和最後的
表，並指出主要群聚和次要群聚。
\item 能說明刪除為什麼要用標記，正確地處理標記上的搜尋和插入。
\item 能在均勻雜湊下推導失敗搜尋 $\frac{m+1}{m-n+1} < \frac{1}{1-\alpha}$ 和成功
搜尋 $\le \frac{1}{\alpha}\ln\frac{1}{1-\alpha}$，記得線性探測的 Knuth 公式，並用它
們算出該把 $\alpha$ 控制在多少。
\item 能追蹤可延伸雜湊的插入，並說明通用雜湊的定義和它保證了什麼。
\item 能做布隆過濾器的插入與查詢，說明它為什麼只有偽陽性、為什麼不能直接刪除，推導偽陽性
率 $(1 - e^{-kn/m})^k$ 和最好的 $k = \frac{m}{n}\ln 2$。
\end{itemize}

\textbf{延伸閱讀}（書末會整理完整的參考資料）

\begin{itemize}
\item Thomas H. Cormen 等，\textit{Introduction to Algorithms}，第 4 版，第 11 章：
兩種方法的分析、通用雜湊族的證明，以及不用刪除標記的線性探測刪除。
\item Donald E. Knuth，\textit{The Art of Computer Programming}，第 3 冊，第 2 版，
第 6.4 節：雜湊的原始整理，包括線性探測公式的完整推導和各種探測法的比較。
\item Ronald Fagin、Jürg Nievergelt、Nicholas Pippenger、H. Raymond Strong，
``Extendible Hashing---A Fast Access Method for Dynamic Files''，\textit{ACM
Transactions on Database Systems} 4(3)，1979：可延伸雜湊的原始論文。
\item J. Lawrence Carter、Mark N. Wegman，``Universal Classes of Hash Functions''，
\textit{Journal of Computer and System Sciences} 18(2)，1979：通用雜湊的原始論文。
\item Steven S. Skiena，\textit{The Algorithm Design Manual}，第 3 版，第 3.7 節與
第 6 章：雜湊在實務上的用法，以及用機率看雜湊的角度。
\item Burton H. Bloom，``Space/Time Trade-offs in Hash Coding with Allowable Errors''，
\textit{Communications of the ACM} 13(7)，1970：布隆過濾器的原始論文。
\item Andrei Broder、Michael Mitzenmacher，``Network Applications of Bloom Filters: A
Survey''，\textit{Internet Mathematics} 1(4)，2004：偽陽性率的推導、計數型布隆過濾器和
各種應用。
\end{itemize}

\addcontentsline{toc}{section}{觀念題解答}
\section*{觀念題解答}

\begin{enumerate}[itemsep=0.15em]
\item (a)、(c)、(d)、(f)。 (b)：遇到同一段群聚就會走同樣的格子，這正是主要群聚。
(d)：奇數和 $16$ 互質（定理~\ref{thm:hash-double}）。 (e)：會切斷經過這一格的探測序
列。 (f)：精確值 $\frac{m+1}{m-n+1} < \frac{1}{1-0.8} = 5$。

\item (a) $6083 = 97 \times 62 + 69$，$h = 69$。 (b) $6083A = 3759.50075\ldots$，小數部
分乘以 $1000$ 是 $500.75\ldots$，$h = 500$。 (c) $6083^2 = 37002889$，第 $4$、$5$ 位
是 $0$、$2$，$h = 2$。 (d) 移位：$526 + 803 + 147 = 1476$，$h = 476$；
邊界：$526 + 308 + 147 = 981$，$h = 981$。

\item (a) 起始位置依序是 $8, 8, 9, 10, 8, 2, 0$，試過的格子依序是：$63$：8；
$85$：8, 9；$75$：9, 10；$54$：10, 0；$96$：8, 9, 10, 0, 1；$24$：2；$33$：0, 1, 2, 3。
最後的表（第 $0$ 到 $10$ 格）：$54, 96, 24, 33, -, -, -, -, 63, 85, 75$。 (b) 成功：$(1+2+2+2+5+1+4)/7 = 17/7 \approx 2.43$。 失敗：起始位置 $0$ 到 $3$ 分別要
$5, 4, 3, 2$ 次，$4$ 到 $7$ 各 $1$ 次，$8, 9, 10$ 分別要 $8, 7, 6$ 次，合計 $39$，平均
$39/11 \approx 3.55$。 (c) 刪除 $54$：試 $10, 0$，第 $0$ 格改成標記。 搜尋 $33$：$0$（標記，繼續）、$1, 2,
3$，找到，$4$ 次。 搜尋 $96$：$8, 9, 10, 0$（標記）、$1$，找到，$5$ 次。 插入 $44$
（$h = 0$）：第 $0$ 格是標記，記下來，繼續看 $1, 2, 3$，第 $4$ 格是空的，停；$44$
放進第 $0$ 格，共探測 $5$ 次。

\item 七個鍵的起始位置都是 $5$。 $i^2 \bmod 11$ 依序是 $0, 1, 4, 9, 5, 3, 3, 5, 9, 4,
1$，所以探測序列是 $5, 6, 9, 3, 10, 8, 8, 10, 3, 9, 6$。 前六個鍵依序放在
$5, 6, 9, 3, 10, 8$。 $71$ 整條序列都是滿的，插入失敗，但還有 $5$ 個空格
（第 $0, 1, 2, 4, 7$ 格）。 由定理~\ref{thm:hash-quad}，
二次探測只看得到 $(11+1)/2 = 6$ 個不同的格子，這六格都被起始位置相同的鍵佔了。

\item $h_2$ 依序是 $6, 8, 7, 5, 9, 3, 4$。 試過的格子：$63$：8；$85$：8, 5；
$75$：9；$54$：10；$96$：8, 6；$24$：2；$33$：0。
最後的表：$33, -, 24, -, -, 85, 96, -, 63, 75, 54$。 共探測 $9$ 次，線性探測是 $17$ 次。

\item (a) $0, 1, 4, 1, 0, 1, 4, 1$，只有 $0, 1, 4$ 三種，起始位置 $0$ 的鍵只試得到第
$0, 1, 4$ 格。 (b) 依序插入三個起始位置都是 $0$ 的鍵（例如 $0, 8, 16$），佔滿第 $0,
1, 4$ 格；再插入 $24$，八次探測都落在這三格，失敗，但表裡還有 $5$ 個空格。 (c)
$i(i+1)/2$ 依序是 $0, 1, 3, 6, 10, 15, 21, 28$，除以 $8$ 的餘數是
$0, 1, 3, 6, 2, 7, 5, 4$，八格都走到了。

\item 假設 $0 \le i < j \le m - 1$ 而兩個位置相同，則 $2^p$ 整除
$\frac{j(j+1)}{2} - \frac{i(i+1)}{2} = \frac{(j - i)(j + i + 1)}{2}$，所以 $2^{p+1}$
整除 $(j - i)(j + i + 1)$。 兩個因數的和 $2j + 1$ 是奇數，所以一奇一偶，$2^{p+1}$
必須整除偶數的那個。 但
$1 \le j - i \le m - 1 < 2^{p+1}$，$1 \le j + i + 1 \le 2m - 2 < 2^{p+1}$，兩個都是比
$2^{p+1}$ 小的正數，不可能被它整除，矛盾。

\item (a) 第 $j$ 格是空的，表示 $100$ 個鍵都沒有落在它，機率是 $(1 - 1/100)^{100}$。
對 $100$ 格的指示變數取期望值再相加，得到 $100 \times 0.99^{100} \approx 36.6$ 個。
(b) 成功：$1 + 99/200 = 1.495$。 失敗：$\alpha = 1$。 (c) 有的串列長 $2$、$3$ 甚至更
長，平均要是 $1$，就要有對應數量的空串列。

\item 均勻雜湊：$\frac{1}{1-\alpha} \le 3$，$\alpha \le 2/3$，約可放 $666$ 個。 線性探
測：$\frac12\bigl(1 + \frac{1}{(1-\alpha)^2}\bigr) \le 3$，$(1-\alpha)^2 \ge 1/5$，
$\alpha \le 1 - 1/\sqrt5 \approx 0.553$，約可放 $552$ 個。 範例~\ref{ex:hash-probe-count}
的實測（線性 $\alpha = 0.5$ 時 $2.49$；雙重 $0.75$ 時 $3.80$）和這兩個界相符。

\item H = \texttt{000}11 查目錄 \texttt{000}，指向 A 的桶，還有空位，放進去：\{A, H\}，
沒有拆桶。 I = \texttt{001}01 查目錄 \texttt{001}，也指向這個桶，滿了；局部深度
$2 < 3$，目錄不變，只拆桶。 依第 $3$ 個位元分：H = 00\texttt{0}11 到 \texttt{000}，
A = 00\texttt{1}10 和 I = 00\texttt{1}01 到 \texttt{001}，兩個桶的局部深度都變成 $3$，
其他桶和目錄不變。 D、F 的桶局部深度 $2$，有 $2^{3-2} = 2$ 個
目錄項指向它。

\item (a) 是。 $|\mathcal{H}|/m = 1.5$，每一對鍵最多只能被一個函數撞。 $h_1$ 撞
$\{a, b\}$、$\{c, d\}$，$h_2$ 撞 $\{a, c\}$、$\{b, d\}$，$h_3$ 撞 $\{a, d\}$、
$\{b, c\}$，六對各一次。 (b) 不是。 $h_3'$ 撞 $\{a, b\}$、$\{a, c\}$、$\{b, c\}$，
$\{a, b\}$ 被 $h_1$、$h_3'$ 撞了 $2$ 次。

\item (a) 位置：$25$：$3, 2$；$47$：$3, 4$；$63$：$8, 5$；$90$：$2, 8$。 最後第
$2, 3, 4, 5, 8$ 格是 $1$：
\[
B[0..10] = 0, 0, 1, 1, 1, 1, 0, 0, 1, 0, 0.
\]
(b) $36$：$36 \bmod 11 = 3$，$\lfloor 36/11 \rfloor = 3$，兩個位置都是第 $3$ 格，是
$1$，回答「可能在」，偽陽性。 $58$：位置 $3, 5$，都是 $1$，「可能在」，偽陽性。 $80$：
位置 $3, 7$，第 $7$ 格是 $0$，「不在」，正確。（$36$ 的兩個位置重合，等於只檢查一個位
元，這是 $h_1$、$h_2$ 不夠獨立的後果。）
(c) 清掉第 $8$、$5$ 格之後，$90$ 的位置是 $2, 8$，第 $8$ 格變成 $0$，回答「不在」，但
$90$ 還在集合裡：偽陰性。 $25$、$47$ 的位置沒有被清掉，仍然答對。
(d) 插入四個鍵之後，第 $2$ 格由 $25$、$90$ 各加一，第 $3$ 格由 $25$、$47$，第 $8$ 格由
$63$、$90$，第 $4$、$5$ 格各一次；刪除 $63$ 時第 $8$、$5$ 格各減一：
\[
\begin{aligned}
\text{插入後：}\ C[0..10] &= 0, 0, 2, 2, 1, 1, 0, 0, 2, 0, 0, \\
\text{刪除 $63$ 後：}\ C[0..10] &= 0, 0, 2, 2, 1, 0, 0, 0, 1, 0, 0.
\end{aligned}
\]
查 $90$：$C[2] = 2$、$C[8] = 1$ 都大於 $0$，「可能
在」，正確。 查 $58$：$C[5] = 0$，「不在」；刪除 $63$ 之後，原本的偽陽性也消失了。

\item (a) $m/n = 12$，$k^{*} = 12 \ln 2 \approx 8.32$。 比較兩個整數：$k = 8$ 時
$(1 - e^{-8/12})^8 = (1 - 0.5134)^8 \approx 0.00314$；$k = 9$ 時
$(1 - e^{-9/12})^9 = (1 - 0.4724)^9 \approx 0.00317$。 取 $k = 8$，偽陽性率約
$0.31\%$，這時 $0$ 位元的比例約 $e^{-8/12} \approx 51\%$，接近定理~\ref{thm:hash-bloom-k}
說的一半。
(b) 用最佳的 $k$，$\varepsilon = 2^{-k^{*}} \le 0.001$，$k^{*} \ge \lg 1000 \approx 9.97$，
每個鍵要 $k^{*}/\ln 2 = \ln 1000 / (\ln 2)^2 \approx 14.4$ 個位元，$n = 100$ 時大約
$1438$ 個位元，$k$ 取 $10$。 驗算：$m = 1438$、$k = 10$，$(1 - e^{-1000/1438})^{10}
\approx 0.000999$，剛好在 $0.1\%$ 以下。
(c) $m$、$n$ 固定時，$k$ 越大，被設成 $1$ 的位元越多。 $k = 8$ 時約 $49\%$ 的位元是 $1$，
$k = 20$ 時 $1 - e^{-20/12} \approx 81\%$ 是 $1$，偽陽性率變成 $0.81^{20} \approx 1.5\%$，
是 $k = 8$ 時 $0.31\%$ 的將近五倍。 檢查的位元多了，但每個位元是 $1$ 的機率也變大，兩
者要取平衡，最好的 $k$ 讓一半的位元是 $0$。
\end{enumerate}

\addcontentsline{toc}{section}{歷屆試題詳解}
\section*{歷屆試題詳解}

\pastsol{px:hash-1}{成大資訊聯招 110 學年度 第 3 題}
題目沒有說「幾次碰撞」怎麼數。 這裡照最常見的算法：依題目給的順序把字串放進表裡，放進去時
起始位置已經有別的字串，就算一次碰撞。 另一種數法是數「落在同一格的字串有幾對」，答案寫在
括號裡。

(1) 四個字串都是 a 到 f 各用一次，只是順序不同，字元碼的和都是
\[
97 + 98 + 99 + 100 + 101 + 102 = 597, \qquad 597 = 11 \times 54 + 3,
\]
四個字串的位置都是 $3$。 第一個放進去沒有碰撞，後面三個都撞上，共 $3$ 次碰撞（落在同一格
的有 $\binom{4}{2} = 6$ 對）。 這就是第~\ref{sec:hash-func}~節說的：只把字元碼加起來，用
同一組字母的字串一定撞在一起。

(2) 第 $i$ 個字元乘上 $i$（$i = 1, \ldots, 6$）再加起來：
\begin{center}
\small
\begin{tabular}{lll}
\toprule
字串 & 加權和 & 除以 $11$ 的餘數 \\
\midrule
abcdef & $1 \cdot 97 + 2 \cdot 98 + 3 \cdot 99 + 4 \cdot 100 + 5 \cdot 101 + 6 \cdot 102 = 2107$ & $2107 - 11 \times 191 = 6$ \\
bcdefa & $1 \cdot 98 + 2 \cdot 99 + 3 \cdot 100 + 4 \cdot 101 + 5 \cdot 102 + 6 \cdot 97 = 2092$ & $2092 - 11 \times 190 = 2$ \\
cdefab & $1 \cdot 99 + 2 \cdot 100 + 3 \cdot 101 + 4 \cdot 102 + 5 \cdot 97 + 6 \cdot 98 = 2083$ & $2083 - 11 \times 189 = 4$ \\
defabc & $1 \cdot 100 + 2 \cdot 101 + 3 \cdot 102 + 4 \cdot 97 + 5 \cdot 98 + 6 \cdot 99 = 2080$ & $2080 - 11 \times 189 = 1$ \\
\bottomrule
\end{tabular}
\end{center}
四個位置 $6, 2, 4, 1$ 都不同，$0$ 次碰撞（$0$ 對）。 乘上位置以後，字元的順序會影響結果，
所以能分開這幾個字串；第~\ref{sec:hash-func}~節用 $B$ 進位（Horner 法）是同樣的想法，而且
權重更分散。

\pastsol{px:hash-2}{政大資科所 115 學年度 第 13 題}
(a) 每個鍵算出 $h_1(k) = k \bmod 20$、$h_2(k) = (k \bmod 7) + 1$，再依序試
$h_1, h_1 + h_2, h_1 + 2h_2, \ldots$（都取 $\bmod 20$）。 $i$ 從 $0$ 開始。
\begin{center}
\small
\begin{tabular}{rcccl}
\toprule
鍵 $k$ & $h_1(k)$ & $k \bmod 7$ & $h_2(k)$ & 依序試過的格子（最後一個是放進去的位置） \\
\midrule
8 & 8 & 1 & 2 & 8 \\
38 & 18 & 3 & 4 & 18 \\
3 & 3 & 3 & 4 & 3 \\
5 & 5 & 5 & 6 & 5 \\
28 & 8 & 0 & 1 & 8（有 8），9 \\
18 & 18 & 4 & 5 & 18（有 38），3（有 3），8（有 8），13 \\
65 & 5 & 2 & 3 & 5（有 5），8（有 8），11 \\
83 & 3 & 6 & 7 & 3（有 3），10 \\
25 & 5 & 4 & 5 & 5（有 5），10（有 83），15 \\
\bottomrule
\end{tabular}
\end{center}
最後的表（沒列出的格子是空的）：
\begin{center}
\small
\begin{tabular}{*{9}{c}}
\toprule
3 & 5 & 8 & 9 & 10 & 11 & 13 & 15 & 18 \\
\midrule
3 & 5 & 8 & 28 & 83 & 65 & 18 & 25 & 38 \\
\bottomrule
\end{tabular}
\end{center}
附帶一提，$M = 20$ 不是質數，而 $h_2(k)$ 可能是 $2, 4, 5, 6$，和 $20$ 不互質；例如 $18$ 的
步長 $5$，探測序列只會經過 $18, 3, 8, 13$ 四格。 由定理~\ref{thm:hash-double}，這種鍵的探測
序列走不遍整張表，表還有空格也可能插不進去。 這一題剛好沒有發生。

(b) $48 \bmod 20 = 8$，$48 \bmod 7 = 6$，$h_2(48) = 7$。 依序看第 $8$ 格（是 $8$，不是 $48$）、
第 $15$ 格（是 $25$）、第 $22 \bmod 20 = 2$ 格（空格），停下來，$48$ 不在表裡。 共
$3$ 次探測（照本章的數法，包括最後看到空格的那一次；如果只數「和鍵比較」的次數，是 $2$ 次）。

(c) 答案是 $n/m$。 這裡把「uniform hash function」解讀成：每個鍵的雜湊值等可能是 $m$ 格中
的任何一格（第~\ref{sec:hash-chain-analysis}~節的簡單均勻雜湊）。 設 $n$ 個鍵裡有 $n_o$ 個
奇數、$n_e$ 個偶數，$n_o + n_e = n$。 固定第 $j$ 格，落在這一格的鍵數是
\[
X_j = \sum_{\text{奇數鍵 } x} I\{h_a(x) = j\} + \sum_{\text{偶數鍵 } x} I\{h_b(x) = j\}.
\]
每個奇數鍵用 $h_a$，落在第 $j$ 格的機率是 $1/m$；每個偶數鍵用 $h_b$，機率也是 $1/m$。 由期
望值的線性，
\[
\mathrm{E}[X_j] = \frac{n_o}{m} + \frac{n_e}{m} = \frac{n}{m} = \alpha .
\]
期望值的線性不需要獨立，所以 $h_a$、$h_b$ 之間有沒有關係都不影響答案；奇偶各用哪個函數也
不影響，只要每個函數本身是均勻的。

\pastsol{px:hash-3}{中央資工所 115 學年度 第 12 題}
題目的 ``liner probing'' 是 linear probing 的筆誤。 題目沒有寫出二次探測的公式，這裡照
第~\ref{sec:hash-open}~節的 $h(k, i) = (h(k) + i^2) \bmod 11$，$i$ 從 $0$ 開始，答案是
A、B、C。 如果改用該節最後提到的正負交替跳法（$h, h + 1^2, h - 1^2, h + 2^2, \ldots$），
A、B 都不成立：$15$ 從第 $4$ 格出發依序試 $4, 5, 3$，在 A、B 裡 $15$ 插入時 $57$ 還沒來，
第 $3$ 格是空的，$15$ 會放在第 $3$ 格，和表不符，答案只剩 C。 作答時先寫明用的是哪一種。

八個鍵的起始位置是：
\[
24 \to 2,\ 57 \to 2,\ 37 \to 4,\ 15 \to 4,\ 27 \to 5,\ 16 \to 5,\ 38 \to 5,\ 17 \to 6 .
\]
判斷方法就是照每個選項的順序真的插一次，看最後的表是不是和題目一樣。 二次探測從 $h$ 出
發依序試 $h, h + 1, h + 4, h + 9, \ldots$。 下面每個鍵寫出試過的格子，最後一個是放進去的
位置。
\begin{itemize}
\item A（二次探測）：$27$：5；$37$：4；$16$：5, 6；$17$：6, 7；$24$：2；$15$：4, 5, 8；
$57$：2, 3；$38$：5, 6, 9。 和題目的表完全相同，正確。
\item B（二次探測）：$37$：4；$27$：5；$16$：5, 6；$24$：2；$15$：4, 5, 8；$57$：2, 3；
$38$：5, 6, 9；$17$：6, 7。 也完全相同，正確。（$17$ 最後才放，第 $7$ 格一直空著，所以
它還是落在第 $7$ 格。）
\item C（線性探測）：$37$：4；$27$：5；$16$：5, 6；$24$：2；$57$：2, 3；$17$：6, 7；
$15$：4, 5, 6, 7, 8；$38$：5, 6, 7, 8, 9。 完全相同，正確。
\item D（線性探測）：$27$：5；$16$：5, 6；$37$：4；$38$：5, 6, 7，放在第 $7$ 格。 但題目的
表第 $7$ 格是 $17$，第 $9$ 格才是 $38$，已經不同，錯誤。 原因是 $38$ 比 $17$ 先插入，搶走了
第 $7$ 格；繼續做下去，最後的表是 $2$：57、$3$：24、$7$：38、$9$：17，和題目不同。
\end{itemize}
也可以只用一個必要條件先刪選項：線性探測裡，鍵 $k$ 放在第 $j$ 格，那麼從 $h(k)$ 到
$j - 1$ 的每一格在 $k$ 插入時都已經有鍵。 D 裡 $38$ 放在第 $9$ 格需要第 $7$ 格先有鍵，也就是
$17$ 要比 $38$ 先插入，但 D 的順序是 $38$ 在前，所以不可能。

\pastsol{px:hash-4}{臺大資工所 115 學年度 第 15 題}
答案是 (C)。
\begin{itemize}
\item (A) 錯。 鏈結法的載入因子可以超過 $1$（第~\ref{sec:hash-chain-c}~節的例子就是
$\alpha = 8/7$），串列只是變長，插入永遠做得到。 「$\alpha = 1$ 就插不進去」是開放定址的情形。
\item (B) 錯。 雜湊表不維持任何順序，再雜湊時每個鍵用新的 $m$ 重算位置
（第~\ref{sec:hash-dynamic}~節），不會排序，之後的搜尋也不用二分搜尋。
\item (C) 對。 在簡單均勻雜湊下，串列的期望長度是 $\alpha = n/m$，搜尋的期望時間是
$\Theta(1 + \alpha)$（定理~\ref{thm:hash-chain-len}）。 $m$ 固定不變時，$\alpha$ 隨 $n$ 線
性成長，搜尋就變成 $\Theta(n)$；$\alpha$ 超過門檻就把表加倍，$\alpha$ 一直維持在常數以下，
搜尋的期望時間是 $O(1)$。 搬動的成本平攤到插入上，每次仍是 $O(1)$。
\item (D) 錯。 表變大只讓碰撞的期望次數變少，不能消除碰撞：任何固定的雜湊函數，鍵空間比
表大時一定有鍵撞在一起（第~\ref{sec:hash-func}~節的鴿籠原理），擴表之後再插入的鍵也一樣會
碰撞。
\item (E) 錯。 擴表配置的是新的串列開頭陣列，串列的節點還是一個一個配置的，不保證放在
相鄰的位置；減少快取失誤也不是擴表的目的。
\end{itemize}

\pastsol{px:hash-5}{成大資訊聯招 113 學年度 第 5 題}
位元陣列有 $15$ 格，所以 $m = 15$。

i) 答案是 (C)。 逐一算三個位置 $f_1(k) = 3k \bmod 15$、$f_2(k) = (2k + 1) \bmod 15$、
$f_3(k) = k^2 \bmod 15$：
\begin{center}
\small
\begin{tabular}{cccccl}
\toprule
$k$ & $f_1$ & $f_2$ & $f_3$ & 三格的位元 & 結論 \\
\midrule
12 & $36 \bmod 15 = 6$ & $25 \bmod 15 = 10$ & $144 \bmod 15 = 9$ & $1, 0, 1$ & 插入會把第 $10$ 格改成 $1$ \\
2 & $6$ & $5$ & $4$ & $1, 1, 1$ & 位置確實是 $4, 5, 6$ \\
15 & $45 \bmod 15 = 0$ & $31 \bmod 15 = 1$ & $225 \bmod 15 = 0$ & $1, 1, 1$ & 查詢回答 yes \\
8 & $24 \bmod 15 = 9$ & $17 \bmod 15 = 2$ & $64 \bmod 15 = 4$ & $1, 0, 1$ & 查詢回答 no \\
\bottomrule
\end{tabular}
\end{center}
\begin{itemize}
\item (A) 錯。 第 $10$ 格原本是 $0$，插入 $12$ 會把它改成 $1$。
\item (B) 錯。 $2$ 的三個位置雖然是 $4, 5, 6$，但布隆過濾器不能用清成 $0$ 的方式刪除：這三
個位元可能也屬於別的鍵，清掉之後查那些鍵會答「不在」，變成偽陰性
（第~\ref{sec:hash-bloom}~節）。
\item (C) 對。 $15$ 的位置是第 $0$、$1$、$0$ 格（$f_1$ 和 $f_3$ 剛好相同），都是 $1$，回答
yes。 至於 $15$ 是不是真的插入過，從過濾器看不出來，可能是偽陽性。
\item (D) 錯。 $8$ 的第二個位置第 $2$ 格是 $0$，回答 no。
\item (E) 錯，因為 (C) 正確。
\end{itemize}

ii) 答案是 (C)。 由定理~\ref{thm:hash-bloom-k}，最佳的雜湊函數個數是
$k^{*} = \frac{m}{n}\ln 2$，這裡 $m = 15$ 個位元、$n = 5$ 個元素，$k^{*} = 3 \ln 2 \approx
2.08$。 (A) 是 $m$，(B) 是 $n$，(D) 把 $m/n$ 錯寫成 $n$。 實際要取整數時比較 $k = 2$ 和
$k = 3$：$(1 - e^{-2/3})^2 \approx 0.237$，$(1 - e^{-1})^3 \approx 0.253$，取 $k = 2$。
